\documentclass[11pt,a4paper]{article}
\usepackage[T1]{fontenc}
\usepackage[utf8]{inputenc}
\usepackage{lmodern}
\usepackage[margin=26mm,headheight=15pt]{geometry}
\usepackage{amsmath,amssymb,amsthm,mathtools}
\usepackage{microtype,booktabs,array,tabularx,longtable}
\usepackage{xcolor,enumitem,fancyhdr,needspace,xurl}
\usepackage{hyperref,bookmark}
\definecolor{ink}{HTML}{20384A}
\definecolor{accent}{HTML}{365E64}
\hypersetup{colorlinks=true,linkcolor=ink,citecolor=accent,urlcolor=accent,
 pdftitle={Surreal Probability and Log-Odds},
 pdfauthor={Research article prepared for Vladimir Reshetnikov},
 pdfsubject={Surreal probability, Bayesian inference, logits, information theory, infinite sampling}}
\urlstyle{same}
\pagestyle{fancy}\fancyhf{}
\fancyhead[L]{\small Surreal Probability and Log-Odds}
\fancyhead[R]{\small Foundations and applications}
\fancyfoot[C]{\thepage}
\renewcommand{\headrulewidth}{0.3pt}
\setlength{\parskip}{3.5pt}
\setlength{\parindent}{1.1em}
\setlength{\emergencystretch}{2.5em}
\widowpenalty=10000\clubpenalty=10000
\allowdisplaybreaks[2]
\setcounter{tocdepth}{2}
\numberwithin{equation}{section}
\newtheorem{theorem}{Theorem}[section]
\newtheorem{proposition}[theorem]{Proposition}
\newtheorem{lemma}[theorem]{Lemma}
\newtheorem{corollary}[theorem]{Corollary}
\theoremstyle{definition}
\newtheorem{definition}[theorem]{Definition}
\newtheorem{example}[theorem]{Example}
\newtheorem{convention}[theorem]{Convention}
\theoremstyle{remark}
\newtheorem{remark}[theorem]{Remark}
\newcommand{\R}{\mathbb R}
\newcommand{\Q}{\mathbb Q}
\newcommand{\N}{\mathbb N}
\newcommand{\Z}{\mathbb Z}
\newcommand{\No}{\mathbf{No}}
\newcommand{\E}{\mathbb E}
\newcommand{\A}{\mathcal A}
\newcommand{\B}{\mathcal B}
\newcommand{\F}{\mathcal F}
\newcommand{\st}{\operatorname{st}}
\newcommand{\lc}{\operatorname{lc}}
\newcommand{\supp}{\operatorname{supp}}
\newcommand{\Var}{\operatorname{Var}}
\newcommand{\Cov}{\operatorname{Cov}}
\newcommand{\logit}{\operatorname{logit}}
\newcommand{\LSE}{\operatorname{LSE}}
\newcommand{\TV}{\operatorname{TV}}
\newcommand{\Fin}{\operatorname{Fin}}
\newcommand{\ind}{\mathbf 1}
\newcommand{\eps}{\varepsilon}
\newcommand{\hsum}{\mathop{\sum}\nolimits^{\mathrm H}}
\newcommand{\abs}[1]{\lvert#1\rvert}
\newcommand{\norm}[1]{\lVert#1\rVert}
\newcommand{\pin}{\texttt{d22a5b35d5b3040e870c3cfd6d1c8f7259e094b0}}
\newcommand{\repo}{\texttt{VladimirReshetnikov/Surreal}}
\newcommand{\smallo}{\mathrm{o}}
\setlist[itemize]{leftmargin=1.5em,itemsep=2pt,topsep=4pt}
\setlist[enumerate]{leftmargin=1.8em,itemsep=3pt,topsep=4pt}

\begin{document}
\hypersetup{pageanchor=false}
\begin{titlepage}
\centering
{\large\sffamily\color{accent} SURREAL MATHEMATICS}\par
\vspace{0.5in}
{\Huge\bfseries\color{ink} Surreal Probability\\[7pt] and Log-Odds}\par
\vspace{0.22in}
{\Large A multiscale theory of belief, information,\\[4pt] and infinite sampling}\par
\vspace{0.24in}
{\large Research article prepared for Vladimir Reshetnikov}\par
\vspace{0.09in}
{22 September 2026}\par
\vspace{0.35in}
\begin{minipage}{0.96\textwidth}
\begin{abstract}
We develop probability with values in set-sized subfields of the surreal
numbers, and logarithmic odds defined by the canonical surreal exponential.
The finite theory includes conditioning on infinitesimal events, exact
Bayesian updates, a leading-scale calculus with coefficient-sensitive tie
breaking, quantitative conditioning stability, softmax models, relative
entropy, proper scoring rules, Gibbs variational principles, and
finite-horizon stochastic processes. A conditional real shadow retains
information that the unconditional standard part destroys. A finite
linear-algebra compression also captures all comparisons between real-valued
acts without discarding signed correction coefficients.

Infinite sample spaces require additional structure rather than a change of
scalar field alone. We distinguish strong Hahn summation, coefficientwise
countable additivity, and regular finitely additive probabilities obtained
from ultrafilters or ordered-field compactness. We construct diffuse
hierarchical probabilities and posterior kernels, prove two elementary
obstructions to coefficientwise independent sampling, and give a regular
all-subsets extension theorem after a set-sized field enlargement. An
uncountable-support example shows that naive coefficientwise integration can
send a pointwise positive surreal-valued function to a negative number.
Standard-part and nonstandard/Loeb bridges are treated separately. Throughout,
finite algebraic identities, support-controlled series, and probabilistic
limit theorems are kept distinct.
\end{abstract}
\end{minipage}
\vfill
\begin{minipage}{0.96\textwidth}
\small
\textbf{Status and provenance.} This is an AI-assisted mathematical synthesis
with detailed proofs and executable exact finite checks, not an independently
refereed or Lean-verified paper. Non-Archimedean probability, lexicographic
belief, and extended log-likelihood methods have substantial precedents.
No first-in-literature claim or solution of a named open conjecture is made.
The repository comparison is pinned to the commit recorded in
Section~\ref{fsp:sec:position}. The article was delivered as a standalone
file and is now a single-source report of that collection;
Section~\ref{fsp:sec:provenance} records its placement and every change.
No statement of it has a Lean formalization.
\end{minipage}
\end{titlepage}
\hypersetup{pageanchor=true}
\pagenumbering{roman}
{\small\setlength{\parskip}{0pt}\tableofcontents}
\clearpage\pagenumbering{arabic}

\section{What a surreal probability theory should accomplish}
\label{fsp:sec:position}

A probability can be smaller than every positive real number without being
zero. Such a value distinguishes an event that is negligible at ordinary
resolution from an event excluded by the model. For example, with
$t=\omega^{-1}$, the probabilities
\begin{equation}\label{fsp:eq:intro}
 (p_1,p_2,p_3)=\frac{(1,t,t^2)}{1+t+t^2}
\end{equation}
encode three different orders of plausibility. Their ordinary shadow is
$(1,0,0)$, but conditioning on $\{2,3\}$ gives
$(1/(1+t),t/(1+t))$ on that event. The shadow alone cannot recover this answer.
A logit can likewise be positive infinite while its probability remains
strictly below one.

The proposal is not to permit probabilities larger than one. Every
probability in this article belongs to $[0,1]$ of an ordered field, hence is a
finite surreal. Infinite quantities arise as logits, likelihood ratios,
surprisals, utilities, or expectations. Nor is the proposal to replace every
ordinary probability calculation: a real probability is already surreal.
The added value is the retention of specified, ordered scales of negligible
mass and evidence.

\subsection{A core theory and several infinite extensions}

There are two logically separate tasks. The first is finite probability
algebra over an ordered field, augmented by a suitable exponential for
information theory. The second is the interpretation of an infinite family
of events, observations, or payoffs. The first task has a broad, consistent
answer. The second has several useful but inequivalent answers.

\begin{center}
\small
\begin{tabularx}{\textwidth}{@{}>{\raggedright\arraybackslash}p{0.20\textwidth}>{\raggedright\arraybackslash}X>{\raggedright\arraybackslash}X@{}}
\toprule
Regime & What is required & Appropriate use\\
\midrule
Finite algebraic & Positivity, normalization, finite additivity & Exact Bayes,
finite decisions, finite graphs and experiments\\
Strong Hahn & Well-ordered union of supports; finitely many contributions per
exponent & Support-controlled infinite algebra and special atomic laws\\
Coefficientwise & A common well-ordered support of ordinary finite signed
measures & Diffuse real laws with infinitesimal layers\\
Regular finite-additive & A chosen extension, often ultrafilter based & Positive
mass for every nonempty event, all-subsets models\\
Nonstandard/Loeb & Internal spaces and suitable saturation, followed by
standard part & Ordinary countably additive limits of hyperfinite experiments\\
\bottomrule
\end{tabularx}
\end{center}

These rows are not five names for the same construction. In particular,
``countably additive surreal probability'' is incomplete terminology unless
the addition and convergence conventions are specified.

\subsection{Relation to the supplied repository}\label{fsp:sub:repository}

The repository was consulted at commit
\begin{center}\small\pin\end{center}
Its \emph{Hahn-Valued Measures and Probability} report already distinguishes
strong and coefficientwise measures, develops strong atomicity and extension
criteria, and studies infinite products and hidden negative mass
\cite{repomeasures}. Its guide explicitly restricts finite-support standard
parts and the support-order threshold $\omega+1$ to the strong class. We do
not generalize those conclusions to other regimes. The foundations report
separates formal Hahn sums from fine or valuation convergence
\cite{repofound}; the Markov report concerns finite generators and scale
resolvents rather than general path measures \cite{repomarkov}.

This article complements those topics by emphasizing conditional shadows,
Bayesian and logit geometry, finite information theory, and the interfaces
between infinite probability constructions. The two nonextension arguments
in Section~\ref{fsp:sec:iid} are elementary special cases consistent with the
repository's broader product discussion; they are not presented as new
independent classification theorems. The repository's maintained guide and
the opening of its source were read for the comparison. This is not a
line-by-line independent audit of all of its proofs or formalization claims.

\paragraph{Since the pin.}
The manuscript was placed in the collection by commit \texttt{5fe7f8d},
seventeen commits after its pin. The statements above remain true of the pin
and are kept as provenance. The comparison of
Section~\ref{fsp:sec:neighbours} was made at commit \texttt{50cb709}, against
the full sources of the three reports cited above and of the physics report
\cite{repophys}. Between the pin and \texttt{50cb709} the source of the
measures report did not change. The placement commit added to that report's
directory the program, build script and records of a further manuscript, on
coefficientwise measure theory, with files prefixed
\texttt{16-coefficientwise-measure-theory-}; this report makes no statement
about that manuscript. The Markov report gained one comparison paragraph about
another report, and the foundations and physics reports are unchanged. At
\texttt{50cb709} a full-text search of every LaTeX source of the collection
finds the words \texttt{logit}, \texttt{softmax}, \texttt{relative entropy},
\texttt{Bayes}, \texttt{martingale} and \texttt{Gibbs} only in this report,
and \texttt{Loeb} elsewhere only in the physics report. Such searches locate
words, not results stated under other names.

\subsection{Precedents and the meaning of development}

The scalar ingredients come from Conway and Gonshor and the exponential-field
results of van den Dries and Ehrlich \cite{conway,gonshor,vdde,vddeerr}.
Set-sized exponential/logarithmic workspaces are also studied explicitly by
Bournez and Guilmant \cite{bournez}. Non-Archimedean probability on infinite
spaces, with a nonclassical infinite-addition rule, is developed by Benci,
Horsten, and Wenmackers \cite{bhw}. Halpern and Brickhill--Horsten analyze
connections with conditional and lexicographic probability
\cite{halpern,brickhill}. Hammond already discusses non-Archimedean
probability representations and extended logarithmic likelihood ratios
\cite{hammond}. Ordinal ranking theories, including Spohn's, are an important
precedent for scale-sensitive belief revision \cite{spohn}. Loeb's theorem
provides a different bridge from nonstandard to real measures \cite{loeb}.

Our development consists of precise compatible definitions, detailed
consequences, counterexamples, and constructions. It does not depend on
asserting historical novelty. Some propositions are useful combinations of
standard facts; others are direct calculations specific to this exposition.
A proof supplied here establishes the stated implication, not priority over
unsearched literature.

\subsection{What this report does not claim}\label{fsp:sec:nonclaims}
Each limitation below is also stated where it applies. They are collected
here, from the article, its delivered README, research audit and verification
record, so that none of them is lost when results are quoted elsewhere.
\begin{enumerate}[label=(N\arabic*),leftmargin=3.2em]
\item \emph{Status and priority.} The report is an AI-assisted synthesis with
written proofs. It has not been independently refereed, no statement of it has
a Lean formalization, and no Lean code accompanies it. Non-Archimedean
probability \cite{bhw}, lexicographic and conditional probability
\cite{halpern,brickhill}, extended logarithmic likelihood ratios
\cite{hammond}, ranking theories \cite{spohn} and Loeb measure \cite{loeb} are
precedents. No first-in-literature result is claimed, no named open
conjecture is claimed solved, and no absence claim about the published
literature is made. Some propositions are combinations of standard facts; a
proof establishes its implication, not priority. A bibliography entry does not
certify every theorem of its source. Loeb's paper was cited from its indexed
metadata, since direct access to the full text failed; the Brickhill--Horsten
citation is to the preprint; Conway and Gonshor are cited as foundational books;
and no ordinal-length estimate of the uncorrected van den Dries--Ehrlich paper is
used. Ordinary real disintegration, product laws and strong laws of large
numbers are imported from \cite{kallenberg}, not re-proved.
\item \emph{Repository comparison.} The comparison at the pin read the guide of
the measures report and the opening of its source; the descriptions of the
foundations and Markov reports rested on the collection's guides. It was not a
line-by-line audit of any report or of the Lean ledger. The strong-class
conclusions of the measures report (finitely supported standard part, the
$\omega+1$ threshold) are not generalized to other regimes. The two coin
obstructions of Section~\ref{fsp:sec:iid} are special cases of that report's
product theorems (Section~\ref{fsp:sec:neighbours}), not new classification
theorems.
\item \emph{Scalars and size.} All sample spaces, event algebras, supports and
parameter collections are set-sized; no probability distribution over the
class $\No$ is constructed. The closure of Lemma~\ref{fsp:lem:workspace}
supplies no algorithm, no canonical smallest
Hahn presentation and no elementary transfer for arbitrary formulas. A Hahn
field $K_\Gamma$ need not be closed under the canonical exponential. The
canonical logarithm is not the omega-map, and $\omega^\gamma=\exp(\gamma\log\omega)$
is not assumed for arbitrary $\gamma$. Transfer from the real exponential field
is used only for universally quantified inequalities in a fixed ordinary number
of variables, never for statements quantifying over sample spaces, sequences,
measurable functions or integration.
\item \emph{The finite core.} Definition~\ref{fsp:def:probability} carries no countable additivity.
Finite Jensen is algebraic convexity, not a statement about a function given
only on $\R$ without a specified extension. Zero-mass cells do not determine
conditional probabilities. The finite no-sure-loss argument does not license
infinitely many tickets, and physical interpretations of infinite stakes are a
separate modeling issue.
\item \emph{Shadows and compression.} Standard part forgets small quantities;
it does not show that their effect after division by other small quantities is
negligible. The lexicographic description after
Theorem~\ref{fsp:thm:skeleton} represents conditional shadows only: it does
not reconstruct the probability or all expected-utility comparisons.
Theorem~\ref{fsp:thm:compression} concerns real-valued acts; it does not say
that an arbitrary signed-row list is a lexicographic probability system, and
it does not compress comparisons of surreal-valued acts.
\item \emph{Bayes and logits.} No unspecified infinite product of likelihoods is
licensed; learning from an infinite data stream needs a chosen law on streams
and a chosen limit. The exponents of the minimum-plus rule lie in an ordered
abelian group, and their addition is not ordinal addition. The endpoints
$p=0,1$ correspond to formal symbols $\mp\infty$ outside $\No$, not to
$\mp\omega$. Logits do not add over disjoint events; a score becomes a
probability only through a normalization rule; stable normalization gives no
finite symbolic representation of an arbitrary surreal score.
\item \emph{Information and optimization.} All sums are finite. The value
$+\infty$ of a relative entropy is formal. Entropy can leave a pure Laurent
field, and no single Hahn field is assumed closed under the information
operations. No entropy rate of a process is supplied, and standard part does
not commute with relative entropy. The logistic optimum comes from an explicit
identity, not an extreme-value theorem. The Gibbs principle constructs no
thermodynamic limit, infinite-volume Gibbs measure or partition function on an
infinite state space. The standard part of an expected utility needs limited
payoffs, and surreal arithmetic does not decide which utilities and trades are
admissible. The smoothing constants and scale are modeling choices, and the
case $N=0$ is an ordinary prior. Proposition~\ref{fsp:prop:separation} gives
no existence theorem for constrained or regularized regression.
\item \emph{Finite processes.} Consistency of finite-horizon path weights is not
a law on infinite paths; a stochastic matrix, a semigroup and a path measure are
different objects. Optional stopping is proved for bounded stopping times only.
``Take $n=1/t$'' has no meaning until a nonstandard or surreal-indexed number of
trials is defined. No unconditional law of large numbers in the fine topology is
hidden in the finite variance bounds.
\item \emph{Infinite addition.} Proposition~\ref{fsp:prop:fine} concerns the full
fine topology and set-sized index sets, not the intrinsic topology of a subfield.
Failures of the strong rule are failures of a support rule, not contradictions
in the field or in finite additivity. The strong classification of the measures
report is quoted, not re-proved. Positivity on cylinders or finite tests does
not imply positivity on all measurable sets.
\item \emph{Hierarchies and integration.} Setting $t=1$ in a hierarchy is not an
ordinary mixture. Dominated convergence holds coefficientwise, not in the fine
or order topology. The Fubini identity is support-controlled and does not cover
all surreal-valued functions. Theorem~\ref{fsp:thm:integrationfailure} asserts
coefficientwise measurability only and does not contradict Lebesgue integration;
the countable-support positivity statement after it is a sufficient condition.
The posterior kernels of Section~\ref{fsp:subsec:kernels} are not a general
surreal Radon--Nikodym or disintegration theorem; their version choices remain
real choices, and no unique conditional law is assigned to observations of
exactly zero mass. Example~\ref{fsp:ex:diffuse} is not an all-subsets regular
probability.
\item \emph{Coin obstructions.} Theorems~\ref{fsp:thm:rarecoin}
and~\ref{fsp:thm:faircoin} do not make any finite coin experiment inconsistent,
do not exclude a regular finitely additive path law
(Corollary~\ref{fsp:cor:coinextension} constructs one), and do not say that every
non-Archimedean Kolmogorov-type theorem fails. A logit parametrization supplies
no infinite-addition rule.
\item \emph{Extensions.} The embeddings of Lemma~\ref{fsp:lem:embedding} need
not preserve an exponential, internal-set structure or a Hahn summation, and do
not show that two exponential structures agree. The ultrafilter construction is
nonconstructive, its snapshot addition is not a strong sum, the resulting laws
are not unique even with a fixed infinitesimal unit, and symmetry is not implied
by equal singleton masses. Theorem~\ref{fsp:thm:extension} provides no countable
additivity, symmetry, canonical embedding or bound on the new scales, and does
not compute the probability of a tail event, which can depend on the model
chosen. Its compactness is first-order logical compactness, not topological
compactness of a simplex.
\item \emph{Latent regime and bridges.} Proposition~\ref{fsp:prop:tailfailure}
is consistent with the real martingale convergence theorem; a surreal analogue
must specify a different conclusion or further hypotheses. Writing an infinite
logit does not define a surreal-indexed sum of log-likelihood increments. The
Loeb bridge needs internal sets and saturation; a field embedding does not
supply them, and a bare surreal $H>\N$ is not a substitute. An internal power
is not identified with a canonical surreal exponential expression under a
field-only embedding.
\item \emph{Not asserted anywhere} (Section~\ref{fsp:sub:notasserted}). No countably
additive probability on all events for the full fine topology; no universal
Carath\'eodory, Radon--Nikodym, Fubini, martingale convergence, strong law of
large numbers, central limit or stochastic calculus theorem after replacing
$\R$ by $\No$; no unique answer to exact-zero conditioning; no canonical
ultrafilter; no reconciliation of regularity with every symmetry demand; no
effective computability of arbitrary surreal expressions; no identification of
the omega-map with exponential powers; no transfer of internal hyperfinite
operations through a field-only embedding. The Poisson and latent examples
import ordinary component laws or a stated nonstandard model.
\item \emph{Computation and evidence.} The engine of
Section~\ref{fsp:sec:implementation} is a specification; a leading-scale answer
is not an exact probability. The program \texttt{code/verify.py} checks finite
examples in $\Q(t)$ with exact rational arithmetic and no floating-point
surrogate for $t$. It implements no canonical exponential or logarithm, no
infinite summation, measure extension, saturation, compactness, choice or
uncountable set, and it proves no general theorem. Compilation and visual
inspection are typesetting checks, not proof checks.
\item \emph{Open.} The four directions of Section~\ref{fsp:sec:scope} are not
settled: positive integration domains for coefficientwise probabilities beyond
bounded real observables and countable supports; convergence of conditional
beliefs that keeps chosen infinitesimal regimes; extension with additional
structure (invariance, a preferred support group, an exponential-compatible
embedding); and a verified computational library.
\end{enumerate}

\subsection{Relation to other reports of the collection}\label{fsp:sec:neighbours}
No proof in this report depends on another report, and no other report
depends on this one. The comparisons below are with the sources as committed
at \texttt{50cb709}; theorem numbers of other reports are those of their
builds at that commit, and each is followed by its label.

\paragraph{Hahn-valued measures and probability.}
That report \cite{repomeasures} is the collection's measure theory for
Hahn-valued probabilities. This report continues it on the finite side and
at its interfaces with other infinite constructions; it answers none of the
questions of Section~33 there.
\begin{enumerate}[label=(\alph*)]
\item \emph{The same definitions.} Definition~\ref{fsp:def:coefficientwise} is
Definition~2.7 there (\nolinkurl{meas:def:coefficientwise}), with the
probability condition of Definition~2.9 there (\nolinkurl{meas:def:positive}).
The strong rule of Section~\ref{fsp:sec:infinity} is Definition~2.8 there
(\nolinkurl{meas:def:strong}), built on the strong summability of Definition~2.2
there (\nolinkurl{meas:def:summability}). Finite additivity
(Definition~\ref{fsp:def:probability} here) is the third notion, and neither report identifies it
with the other two. On a finite event algebra a disjoint sequence has only
finitely many nonempty members, so all three notions coincide there; this is
why Sections~\ref{fsp:sec:shadows} and~\ref{fsp:sec:bayes} may speak of a
\emph{finite Hahn probability} without a qualifier.
\item \emph{Strong class.} Example~\ref{fsp:ex:geometric} is an instance of the
atomic classification, Theorem~4.2 there (\nolinkurl{meas:thm:atomic}). The
``classification'' quoted in Section~\ref{fsp:sec:infinity} is that theorem with
Corollaries~4.5 and~4.6 there (\nolinkurl{meas:cor:two-axioms},
\nolinkurl{meas:cor:shadow}): on countably separated spaces strong measures are
the coefficientwise ones with finitely atomic coefficients, and a strong
probability has a finitely supported standard part. The Lebesgue remark of
Section~\ref{fsp:sec:infinity} is Trap~2 of Section~2.6 there
(\nolinkurl{meas:sub:traps}).
\item \emph{The coin obstructions.} Theorem~\ref{fsp:thm:rarecoin} (bias $t$) and
Theorem~\ref{fsp:thm:faircoin} (bias $1/2+t$) are two cases of
Corollary~23.1 there (\nolinkurl{meas:cor:iid}): the i.i.d.\ Bernoulli law with
bias $p$ has a coefficientwise \emph{signed} extension if and only if $p$ is
real. That corollary rests on the total-variation extension test, Theorem~19.1
there (\nolinkurl{meas:thm:extension}). Its proof uses the same two mechanisms:
the first infinitesimal coefficient of bias $t^\alpha$ has level-$n$ variation
$2n$, where Theorem~\ref{fsp:thm:rarecoin} uses the first-success events, and
for an interior bias the variation grows like $\sqrt n$ by H\"older's inequality
and fourth-power expectations, Lemma~20.3 there (\nolinkurl{meas:lem:L1lower}),
as in the proof of Theorem~\ref{fsp:thm:faircoin}. Neither coin law has a strong
extension either (Corollary~10.3 and Example~10.4 there,
\nolinkurl{meas:cor:bernoulli}, \nolinkurl{meas:ex:bernoulli}).
Corollary~\ref{fsp:cor:coinextension} adds what that report leaves aside: its
non-claims 10.1 and 15.1 (Section~32.2 there, \nolinkurl{meas:sub:nonclaims})
say that its negative results do not exclude finitely additive or ultrafilter
probability, and Corollary~\ref{fsp:cor:coinextension} gives both coin models
regular finitely additive laws on all subsets.
\item \emph{Conditioning.} Proposition~24.2 there
(\nolinkurl{meas:prop:conditioning}) shows that conditioning a coefficientwise
probability on an event of positive mass stays in the class, and Example~24.3
there (\nolinkurl{meas:ex:conditioning}) computes a conditional shadow with
infinitely many atoms. Theorem~\ref{fsp:thm:hierarchy} computes every
conditional shadow of a normalized hierarchy, and Section~\ref{fsp:subsec:kernels}
conditions on continuous observations; neither result is in that report.
\item \emph{Observables.} The expectation \eqref{fsp:eq:hierintegral} is the
restriction to normalized hierarchies of the integral of bounded ordinary
observables in Section~24 there (\nolinkurl{meas:eq:coef-integral}, defined for
every coefficientwise measure), with the same coefficientwise dominated
convergence. That report says that integrating arbitrary Hahn-valued functions
would need further joint support and integrability conditions and asserts
nothing about it. Theorem~\ref{fsp:thm:integrationfailure} shows that a common
well-ordered support with bounded integrable coefficients is not enough at
uncountable support rank, and the paragraph after it shows that a countable
support is enough for positivity. Its mechanism, an uncountable union of null
sets, is reminiscent of the null-ideal criterion, Theorem~19.2 there
(\nolinkurl{meas:thm:nullideal}), in which positivity of a coefficientwise
measure is decided on the null ideal common to all earlier coefficients, and
this ideal is that of one control measure when the predecessor set is countable.
The two results concern different objects (integrals of Hahn-valued functions
against an ordinary measure here, event masses of a Hahn-valued measure there),
and neither is derived from the other.
\item \emph{Open questions.} Question~33.9 there (\nolinkurl{meas:q:markov}),
dependence governed by a transition kernel, stays open: Section~\ref{fsp:sec:processes}
has finite horizons only, and the latent model of Section~\ref{fsp:sec:latent} is an
exchangeable mixture of two product laws.
\end{enumerate}

\paragraph{Foundations.}
The strong summability used here is (10.1) of the foundations report
\cite{repofound} (\nolinkurl{found:eq:summability}), and the positive-support
Neumann lemma is its Section~10.2 (\nolinkurl{found:sub:positivesupport}).
Proposition~\ref{fsp:prop:fine} is clause~(2) of Theorem~12.1 there
(\nolinkurl{found:thm:discrete}), with the same proof; that theorem also makes
set-sized subsets discrete and small-index Cauchy nets eventually constant.
The distinctions drawn in Section~\ref{fsp:sec:infinity} between the full fine
topology, the intrinsic topology of a subfield and coefficientwise real
convergence are those of its Section~12 (\nolinkurl{found:sec:topology}) and of
the topology table in \nolinkurl{docs/NOTATION.md}.

\paragraph{Markov generators at every scale.}
The Markov report \cite{repomarkov} studies the normalized resolvent
$R_L(s)=s(sI+L)^{-1}$ of a finite rate matrix over a Hahn field with divisible
value group. Its non-claim~(N1) (Section~1.4 there,
\nolinkurl{markov:sec:nonclaims}) excludes path measures, sample paths and
surreal-valued probabilities on path space. Section~\ref{fsp:sec:processes}
builds finite-horizon path weights from any finite stochastic matrix over $K$.
The following link is an observation of the placement, not of either source.
By Proposition~3.1 there (\nolinkurl{markov:prop:forest}), for positive $s$ the
matrix $R_L(s)$ is entrywise positive and row stochastic with limited entries.
So it may serve as $T$ in Section~\ref{fsp:sec:processes}, and there the path
shadow is the ordinary finite-horizon chain whose transition matrix is the
residue matrix $K_\alpha(c)$ of $R_L(ct^\alpha)$, which that report calls the
shadow at scale $\alpha$. This gives laws of finite paths only; neither report
constructs a law on infinite paths. The relative-stability Theorem~8.1 there
(\nolinkurl{markov:thm:stability}) concerns resolvent entries under relative rate
errors, and Theorem~\ref{fsp:thm:stability} concerns conditional probabilities
under absolute or valuation-controlled errors; they are statements about
different maps.

\paragraph{Surreal scalars and spacetime.}
The probability subsection of the physics report \cite{repophys}
(Section~15.4 there, \nolinkurl{phys:sub:probability}) observes that equal
positive infinitesimal weights on infinitely many outcomes are not strongly
summable (Example~4.3 there, \nolinkurl{phys:ex:summable}); that a finite
independent-trial model with a fixed ordinary number of trials keeps an
infinitesimal event infinitesimally likely; and that non-Archimedean theories
\cite{bhw} and Loeb's construction \cite{loeb} need additional structure and an
account of how records relate to field-valued weights. Section~\ref{fsp:sec:infinity}
repeats the first point, \eqref{fsp:eq:hitting} is the second, the
fine-ultrafilter law of Section~\ref{fsp:sec:extension} gives equal positive
singleton masses by snapshot addition, not by a strong sum, and
Section~\ref{fsp:sec:bridges} states what Loeb's construction needs. Like that
report, this one gives no physical interpretation of field-valued probabilities.

\subsection{Words and symbols used differently elsewhere in the collection}
\label{fsp:sec:words}
Several words and letters of this report mean other things elsewhere. Each has
one meaning here.
\begin{itemize}
\item \emph{Additivity.} \emph{Finite additivity} is Definition~\ref{fsp:def:probability}.
\emph{Strong} additivity is the rule of Section~\ref{fsp:sec:infinity}, and
\emph{coefficientwise} additivity is Definition~\ref{fsp:def:coefficientwise}; both
are the measures report's classes. The snapshot addition of an ultrapower
(Section~\ref{fsp:sec:extension}) and Loeb's countable additivity of a real
standard part (Section~\ref{fsp:sec:bridges}) are further, different rules.
\emph{Strong} always refers to strong Hahn summability; the classical strong law
of large numbers is always written in full.
\item \emph{Regular} is Brickhill and Horsten's regularity: every nonempty event
of the algebra has positive probability. It is not Radon regularity, and the
measures report gives the word no meaning of its own.
\item \emph{Limited} means bounded in magnitude by an ordinary integer, which
\nolinkurl{docs/NOTATION.md} calls \emph{finite}. Here \emph{finite} counts
things (finitely many states, terms, trials or tickets); an ``ordinary finite
$n$'' is a standard natural number, and the ``finite surreal'' of
Section~\ref{fsp:sec:position} is a limited one.
\item \emph{Shadow.} $P_0=\st P$, the conditional shadow $C(A\mid B)$ and the
path shadow are standard parts, $\st$ being the collection's scalar standard
part. The Markov report's shadow is the same operation applied to a resolvent
matrix. No standard-part topology is used.
\item \emph{Fine.} The fine topology is the collection's full surreal fine
topology. A \emph{fine} ultrafilter (Sections~\ref{fsp:sub:ultrafilter} and~\ref{fsp:sub:nonunique}) is one containing
every cone $\{s:F\subseteq s\}$; that sense, from the non-Archimedean probability
literature, has nothing to do with the topology.
\item \emph{Hierarchy} means a normalized hierarchy of ordinary probability laws
(Theorem~\ref{fsp:thm:hierarchy}), not the Markov report's hierarchy of
stochastic projections.
\item \emph{Moment.} The fourth moment $\E S_n^4$ in the proof of
Theorem~\ref{fsp:thm:faircoin} is the expectation of a power of an ordinary random
variable, which the measures report calls a fourth-power expectation; no moment
problem appears here.
\item \emph{Positive}, for a Hahn series, means a positive leading coefficient.
A positive coefficientwise probability is nonnegative on every event, as in the
measures report; positivity of every event is regularity.
\item \emph{Local letters.} $H$ is an entropy $H(p)$ (Section~\ref{fsp:sec:information}),
a hypothesis with data $D$ in \eqref{fsp:eq:logitbayes}, a hyperinteger sample size
(Sections~\ref{fsp:sub:ultrafilter}, \ref{fsp:sub:nonunique} and~\ref{fsp:sub:poisson}) and the latent label of Section~\ref{fsp:sec:latent}.
$E$ is an exponential--logarithmic workspace, and $E_i\in E$ are energies in
Theorem~\ref{fsp:thm:gibbs}; the blackboard $\E$ is expectation. $\lambda_\gamma$
is $-\log t^\gamma$, while $\lambda$ alone is a precision in
Theorem~\ref{fsp:thm:stability}, Lebesgue probability in
Example~\ref{fsp:ex:diffuse} and a Poisson parameter in
Section~\ref{fsp:sec:bridges}. $D$ is a relative entropy $D(p\Vert q)$, a leading
constant in Theorem~\ref{fsp:thm:scalebayes} and a density $D(y)$ in
\eqref{fsp:eq:kerneldensity}. $T$ is a stochastic matrix in
Theorem~\ref{fsp:thm:dataprocessing} and Section~\ref{fsp:sub:paths} and a stopping time in
Section~\ref{fsp:sub:stopping}. $K$ is the scalar field, $K_j$ are conditional kernels and $K_n$ are
counts in Section~\ref{fsp:sec:latent}. $C(A\mid B)$ is a conditional shadow and
$C$ a real constant in Section~\ref{fsp:sec:bayes}. In examples
$t=\omega^{-1}$.
\item The valuation is $v(x)=\min\supp x$, $\lc$ is the leading coefficient, and
the surreal embedding is $t^\gamma\mapsto\omega^{-\gamma}$, as elsewhere in the
collection. $\exp$ and $\log$ are the canonical ordered surreal exponential and
its inverse, never a global surcomplex exponential.
\end{itemize}

\subsection{Provenance of this report}\label{fsp:sec:provenance}
This report is built from one manuscript, \emph{Surreal Probability and
Log-Odds: A multiscale theory of belief, information, and infinite sampling},
dated 22 September 2026. It arrived as the archive
\nolinkurl{surreal_probability.zip}, number~05 of nine archives committed on
arrival at \texttt{2765c8f}. The archive held the LaTeX source
\texttt{surreal\_probability.tex}, its PDF, a README, the research audit
\texttt{RESEARCH\_AUDIT.md}, the program \texttt{verify.py} with its recorded
output \texttt{verification.json}, a build record \texttt{BUILD\_VALIDATION.json}
and a build script. It was placed at commit \texttt{5fe7f8d}. Its repository
comparison is pinned to \texttt{\pin} (Section~\ref{fsp:sub:repository}).

There is no second source, so there was no merge and no choice between sources.
Placement retained every theorem, proof, example and disclaimer of the
manuscript. The manuscript continues the measures report but answers no question it names, so
it was placed as a report of its own rather than as a part of that one.

Placement changed the following and nothing else.
\begin{enumerate}[label=(\alph*)]
\item The source was renamed \texttt{article.tex}. The program and build script
were moved to \texttt{code/}, the two JSON records to \texttt{data/}; they and
\texttt{RESEARCH\_AUDIT.md} are byte-identical to the delivered files. The
delivered PDF and README are not shipped.
\item All 100 labels of the manuscript received the prefix \texttt{fsp:}, and
every reference was updated. Fourteen labels were added, for the new
subsections and for the definition, subsections and appendix that the new text
cites.
\item Sections~\ref{fsp:sec:nonclaims}--\ref{fsp:sec:provenance}, the paragraph
``Since the pin'' in Section~\ref{fsp:sub:repository} and pointer sentences in
Sections~\ref{fsp:sec:infinity} and~\ref{fsp:sec:iid} were added. They contain no
numbered statement or equation, so every theorem, equation and section number of
the manuscript is unchanged. The title-page status paragraph now states the draft
status and the placement.
\item To keep one meaning of \emph{strong}, the title of
Example~\ref{fsp:ex:geometric} (``A regular strong law'') now reads ``A regular
strong Hahn probability'', and the four mentions of classical strong laws
(Sections~\ref{fsp:sec:latent} and~\ref{fsp:sub:notasserted}, Appendix~\ref{fsp:app:ledger}) now read
``strong law(s) of large numbers''.
\item Section~\ref{fsp:sub:verifier} and Appendix~\ref{fsp:app:build} now describe the files as placed
and record a rerun of the checks; the bibliography entries of the repository
reports carry a placement note, and an entry for the physics report was added.
\end{enumerate}
Placement left the mathematics unchanged. The subsequent proof review of
Sections~\ref{fsp:sec:scalars}--\ref{fsp:sec:finite} expands the elementary
arguments and supplies the full-event-algebra hypothesis for the point-weight
representation. On a smaller finite algebra it is the measurable atoms, not
necessarily the points, that carry the uniquely determined weights.
It also makes the limited-input and conditional-version conventions explicit.
The live review record is \path{docs/REVIEW.md}; the delivered audit, programs
and data remain unchanged.

\section{Scalars, size, and the two meanings of exponentiation}
\label{fsp:sec:scalars}

\subsection{Set-sized workspaces inside the surreals}

The surreal numbers $\No$ form a proper class, not an ordinary set-valued
sample range. We use set-sized sample spaces, event algebras, supports, and
collections of parameters. The algebraic field $K$ always denotes a
set-sized ordered subfield of $\No$ containing the displayed copy of $\R$.
When logarithms are used we work in a set-sized subfield $E\subseteq\No$
closed under the canonical surreal $\exp$ and under $\log$ on positive
arguments. In a class theory these are sets inside the ambient class; in a
universe formulation all sizes are interpreted relative to the chosen
universes. No probability distribution over the entire proper class $\No$
is constructed.

\begin{lemma}[Local exponential workspace]\label{fsp:lem:workspace}
Every set $S\subseteq\No$ is contained in a set-sized subfield $E$ containing
$\R$ and closed under canonical exponential and positive logarithm.
\end{lemma}
\begin{proof}
Let $E_0$ be the field generated by $\R\cup S$. Given $E_n$, let $E_{n+1}$
be the field generated by
$E_n\cup\exp(E_n)\cup\log((E_n)_{>0})$.
Replacement and finite field expressions make each stage a set.
The union $E=\bigcup_{n<\omega}E_n$ has the required properties because each
input and each finite list of inputs already lies at one stage.
More explicitly, the stages form an increasing chain, so two elements of
$E$ lie together in some $E_n$. Their sum, product and any defined quotient
lie in that field. For $x\in E_n$, its exponential and, when $x>0$, its
logarithm lie in $E_{n+1}$. This proves each claimed closure separately.
\end{proof}

This lemma supplies closure, not an algorithm, a canonical smallest Hahn
presentation, or elementary transfer for arbitrary quantified formulas.
The particular scalar inequalities used below hold in $\No$ and therefore
for their inputs in $E$.

\subsection{Standard part}

An element $x$ is \emph{limited} if $|x|<n$ for some ordinary positive
integer $n$, and \emph{infinitesimal} if $|x|<1/n$ for every such $n$.
Write $x\simeq y$ when $x-y$ is infinitesimal. Every limited surreal has a
unique real standard part, denoted $\st x$, with $x\simeq\st x$.
Indeed, if $|x|<n$, the set $\{a\in\R:a<x\}$ is nonempty and
bounded above in $\R$. Let $r$ be its real supremum. For each real
$\delta>0$, the supremum property gives $r-\delta<x$, while
$r+\delta/2$ exceeds the supremum and hence cannot be strictly below $x$.
Hence $r-\delta<x<r+\delta$. This proves $x\simeq r$.
Two such real numbers differ by a real infinitesimal and therefore agree.
The Dedekind completeness used here belongs to $\R$, not to $\No$.
Here \emph{limited} means what the collection's notation guide calls a
\emph{finite element}; it is unrelated to finite birthday or finite support.

\begin{proposition}[Residue algebra]\label{fsp:prop:st}
On limited elements, standard part is an order-preserving ring homomorphism.
If $x,y$ are limited and $\st y\ne0$, then $x/y$ is limited and
$\st(x/y)=\st x/\st y$. A limited $x$ is infinitesimal exactly
when $\st x=0$.
\end{proposition}
\begin{proof}
Write $x=r+\eta$ and $y=s+\theta$ with $r,s\in\R$ and infinitesimal
$\eta,\theta$. Sums and products of infinitesimals with limited factors
are infinitesimal. If $s\ne0$, then $y$ is bounded away from zero by a
positive real, and $y^{-1}-s^{-1}=-\theta/(sy)$ is infinitesimal.
Thus $y^{-1}$ is limited with standard part $s^{-1}$, proving the
quotient assertion. If $x\le y$ but $r>s$, the positive real gap $r-s$
cannot be canceled by the infinitesimal $\theta-\eta$, a contradiction.
Finally $x=r+\eta$ is infinitesimal if and only if $r=0$.
The map fixes real constants, so it is onto $\R$; its kernel is precisely
the infinitesimal ideal of the limited-element ring. It is not strictly
order preserving: $0<t=\omega^{-1}$ but $\st t=\st0=0$.
\end{proof}

A field need not contain all limited surreals to have this residue map;
containing $\R$ is sufficient for the standard parts of its own elements.
Taking standard part forgets small quantities; it does not prove that their
statistical or decision-theoretic effects are negligible after division by
other small quantities.

\subsection{Hahn coordinates and leading scales}

For explicit expansions let $\Gamma$ be a set-sized ordered additive subgroup
of $\No$, and put
\begin{equation}
 K_\Gamma=\R((t^\Gamma)),\qquad t^\gamma\longmapsto\omega^{-\gamma}.
\end{equation}
The exponent support is well ordered in increasing order. These normal
forms identify $K_\Gamma$ with an ordered subfield of $\No$
\cite[Section~2, p.~176]{vdde}; the field need
not be closed under the canonical exponential. An abstract ordered group
can instead be used with a specified embedding into the additive surreals.
For $x\ne0$ set
\begin{equation}
 v(x)=\min\supp(x),\qquad \lc(x)=[t^{v(x)}]x,
 \qquad v(0)=+\infty.
\end{equation}
Here $+\infty$ is a bookkeeping symbol for the valuation, not a surreal.
The order of a nonzero series is the sign of its leading coefficient.
For positive $x,y$,
\begin{equation}\label{fsp:eq:positiveval}
 v(xy)=v(x)+v(y),\qquad
 v(x+y)=\min\{v(x),v(y)\}.
\end{equation}
For products the least exponent is $v(x)+v(y)$, with coefficient
$\lc(x)\lc(y)\ne0$. For sums with unequal valuations, the smaller
exponent survives; with equal valuations its coefficient is
$\lc(x)+\lc(y)>0$. This proves both identities without a cancellation
assumption on any higher coefficient. If $v(p)<0$ and $p>0$, its leading
monomial dominates every positive real constant, contradicting $p\le1$.
Thus probabilities have $v(p)\ge0$. More generally a Hahn element is limited
exactly when its valuation is nonnegative, and its standard part is then
$[t^0]x$; zero has the assigned valuation $+\infty$.

A family $(x_i)_{i\in I}$ is \emph{strongly Hahn summable} if its combined
support is well ordered and, for each exponent, only finitely many members
have a nonzero coefficient there. Its sum, written $\hsum_i x_i$, is formed
coefficientwise using those finite sums. The standard positive-support
Neumann lemma justifies geometric, exponential, and logarithmic formal
series in infinitesimal Hahn arguments; it is a support theorem, not a
statement about ordinary sequence limits \cite{gonshor,repofound}.

\subsection{Canonical logarithm is not the omega-map}

The canonical surreal exponential is an order-preserving group isomorphism
from $(\No,+)$ to $(\No_{>0},\cdot)$ extending real exponential
\cite{gonshor,vdde}. Its inverse is our only meaning of $\log$.
The omega-map $\gamma\mapsto\omega^\gamma$ that defines Conway monomials
is a different operation. In particular, one must not substitute
\begin{equation*}
 \omega^\gamma=\exp(\gamma\log\omega)
 \quad\text{for arbitrary surreal }\gamma
\end{equation*}
without an additional theorem on the exponents in question.
A safe log-scale notation is
\begin{equation}\label{fsp:eq:lambda}
 \lambda_\gamma=-\log(t^\gamma).
\end{equation}
It satisfies $\lambda_{\gamma+\delta}=\lambda_\gamma+\lambda_\delta$ and is
strictly increasing. Thus, for $c>0$ real and infinitesimal $u$,
\begin{equation}\label{fsp:eq:logleading}
 \log\bigl(c t^\gamma(1+u)\bigr)
 =\log c-\lambda_\gamma+\log(1+u).
\end{equation}
Integer powers of a fixed $t$ are ordinary field powers and cause no such
ambiguity.

We use the scalar inequalities
\begin{equation}\label{fsp:eq:logineq}
 \exp x\ge1+x,\qquad \log y\le y-1\quad(y>0),
\end{equation}
with equality only at $x=0$ and $y=1$, respectively. These follow from the
elementary-extension theorem \cite[Corollary~2.2, p.~177]{vdde}.
The erratum \cite{vddeerr} corrects ordinal support bounds in Section~4
and leaves this transfer statement unchanged. No such ordinal bound is
used here. For each fixed ordinary integer $n$, a universally quantified
real-exponential inequality involving $n$ coordinates can be used this way:
first it holds in $\No$, then its quantifier-free instance at inputs in
$E$ has the same meaning there, since $E$ inherits the order and field
operations and is closed under the displayed exponential and logarithm.
The workspace lemma alone does not make $E$ an elementary substructure. This does
\emph{not} transfer a theorem quantifying over
arbitrary sample spaces, countable sequences, measurable functions, or
integration operations. The distinction will matter repeatedly.

\section{The finite probability core}
\label{fsp:sec:finite}

\subsection{Events, weights, and random variables}

\begin{definition}\label{fsp:def:probability}
A $K$-probability on an algebra $\A$ of subsets of a nonempty set $\Omega$
is a function $P:\A\to K$ such that $P(\Omega)=1$, $P(A)\ge0$, and
\begin{equation}
 P(A\cup B)=P(A)+P(B)\quad\text{if }A\cap B=\varnothing.
\end{equation}
It is \emph{regular} if $P(A)>0$ for every nonempty $A\in\A$.
No countable-additivity assertion is implicit in this definition.
\end{definition}

Taking the two disjoint events to be empty gives $P(\varnothing)=0$.
The decompositions $\Omega=A\sqcup A^c$ and
$B=A\sqcup(B\setminus A)$ when $A\subseteq B$ give complements and
monotonicity, hence $0\le P(A)\le1$.

On $\Omega=\{1,\ldots,n\}$ \emph{with the full event algebra
$\A=2^\Omega$}, a law is equivalent to its point-weight vector
$p_i=P(\{i\})\ge0$, with $\sum_i p_i=1$; finite additivity gives
$P(A)=\sum_{i\in A}p_i$, and this formula conversely defines a law.
Regularity means exactly that all $p_i>0$. For a smaller finite event
algebra, use its nonempty minimal events (its \emph{atoms}) instead:
these partition $\Omega$, every event is a union of them, and their masses
uniquely determine the law. Individual point masses need not be defined.
For example, on $\{1,2\}$ with algebra $\{\varnothing,\Omega\}$, the
point-weight vectors $(1,0)$ and $(0,1)$ restrict to the same law.
Throughout the finite vector models below, all subsets are measurable;
this is an additional convention, not a consequence of finiteness alone.

A finite-range $K$-valued random variable is a function constant on the
members of a finite measurable partition. Its expectation is
\begin{equation}\label{fsp:eq:expectation}
 \E_P X=\sum_j x_jP(A_j),\qquad X=\sum_jx_j\ind_{A_j}.
\end{equation}
For two representations, pass to the common partition of nonempty
intersections of their cells. On each such intersection the two assigned
values agree, and finite additivity splits each original mass into the
masses of these intersections. Thus the displayed sum is independent of
the representation. A common partition also proves $K$-linearity and
positivity. Applying positivity to $X-\min X$ and $\max X-X$, and using
$\E1=1$, gives $\min X\le\E X\le\max X$.
No convergence notion is needed, even if the finitely many values are
infinite surreals. If $P(B)=0$, then $\E[X\ind_B]=0$: every cell inside
$B$ has mass zero, whatever its assigned value.

For two finite laws, the product weights $p_iq_j$ define independent
coordinates. A finite map $f:\Omega\to S$ pushes the law forward by
$P_f(s)=\sum_{f(i)=s}p_i$. Every ordinary finite marginalization, finite
inclusion--exclusion identity, and finite factorization argument therefore
works verbatim as field algebra. A finite acyclic graphical model with
stochastic local tables defines a normalized joint law by multiplying
those tables; normalization follows by summing terminal vertices in a
reverse topological order. This is an actual construction, not an appeal
to a missing infinite extension theorem.

\subsection{Finite inequalities}

\begin{theorem}[Variance, Cauchy--Schwarz, and tail bounds]
\label{fsp:thm:finitebounds}
For finite-range measurable variables $X,Y$ on the same $K$-probability
space, define $\Var(X)=\E[(X-\E X)^2]$. Then
\begin{align}
 \Var(X)&=\E[(X-\E X)^2]\ge0,\\
 (\E[XY])^2&\le\E[X^2]\,\E[Y^2],\label{fsp:eq:cs}\\
 P(X\ge a)&\le\E[X]/a\quad(X\ge0,\ a>0),\label{fsp:eq:markov}\\
 P(|X-\E X|\ge a)&\le\Var(X)/a^2\quad(a>0).
\end{align}
The thresholds and variable values may be surreal.
\end{theorem}
\begin{proof}
Work on a common finite measurable partition for $X$ and $Y$.
Positivity of squares gives the variance assertion, and expansion gives
$\Var(X)=\E[X^2]-(\E X)^2$. Put
$A=\E[X^2]$, $B=\E[XY]$, and $C=\E[Y^2]$.
If $C=0$, each nonnegative summand $p_jy_j^2$ is zero, so every
positive-mass cell has $y_j=0$ and $B=0$.
If $C>0$, expansion at $c=B/C$ gives
\[
 0\le\E[(X-cY)^2]=A-2cB+c^2C=A-B^2/C.
\]
Multiplication by $C$ proves \eqref{fsp:eq:cs}.
For Markov, $a\ind_{\{X\ge a\}}\le X$ pointwise; apply positivity
of expectation and divide by $a>0$.
For Chebyshev, apply Markov to $(X-\E X)^2$ at the positive threshold
$a^2$. The events agree because $|z|\ge a$ if and only if $z^2\ge a^2$
in an ordered field. No division by a possibly zero second moment or
probability is used.
\end{proof}

If a function $f$ is convex on an interval of $K$ in the ordered-field
sense, finite Jensen follows by induction on the number of positive
weights. With at least two such weights, separate one of weight $p\in(0,1)$,
normalize the remaining weights by $1-p$, apply the inductive assertion,
and then the defining two-point convexity inequality. One positive weight
is the equality case; zero weights may be discarded. One cannot apply this
assertion to an arbitrary function
originally defined only on $\R$ without first specifying its extension to
$K$. The result is algebraic convexity, not a general surreal integration
theorem.

\subsection{Conditioning and finite conditional expectation}

Whenever $P(B)>0$, define
\begin{equation}\label{fsp:eq:condition}
 P(A\mid B)=\frac{P(A\cap B)}{P(B)}.
\end{equation}
An infinitesimal denominator is nonzero and poses no algebraic problem.
The map $A\mapsto P(A\mid B)$ is again a finitely additive law.
For $P(A),P(B)>0$, cancellation gives
$P(A\mid B)P(B)=P(B\mid A)P(A)$; division by $P(B)$ gives Bayes' rule.
Likewise
$P(A\cap B\mid C)=P(A\mid B\cap C)P(B\mid C)$ whenever
$P(B\cap C)>0$ (which also ensures $P(C)>0$).
Each iterated chain rule requires the displayed conditioning masses to be
positive; an exact zero does not define a quotient.
For a finite partition $(B_j)$ with positive masses,
\begin{equation}
 P(A)=\sum_jP(A\mid B_j)P(B_j).
\end{equation}
Zero-mass cells, when present, contribute zero and may be omitted; assigning
arbitrary conditional values there does not determine a unique conditional
probability system on those cells.

For a finite subalgebra $\mathcal G$ with positive-mass atoms, set
$\E[X\mid\mathcal G]$ on an atom $B$ equal to
$\E[X\ind_B]/P(B)$. Direct finite summation gives
\begin{equation}\label{fsp:eq:tower}
 \E\bigl[\E[X\mid\mathcal G]\bigr]=\E X,\qquad
 \E\bigl[\E[X\mid\mathcal G]\mid\mathcal H\bigr]
 =\E[X\mid\mathcal H]\quad(\mathcal H\subseteq\mathcal G).
\end{equation}
For the first identity, summing over the atoms cancels their positive
masses. For the second, let $C$ be an atom of $\mathcal H$ and sum over
the atoms $B$ of $\mathcal G$ contained in $C$:
\[
 \sum_{B\subseteq C}\frac{\E[X\ind_B]}{P(B)}P(B)
 =\E[X\ind_C].
\]
Dividing by $P(C)$ gives the equality on $C$.
A $\mathcal G$-measurable variable $Z$ is constant on each atom, so
$\E[ZX\mid\mathcal G]=Z\E[X\mid\mathcal G]$ there.

Write $M=\E[X\mid\mathcal G]$ and define
$\Var(X\mid\mathcal G)=\E[(X-M)^2\mid\mathcal G]$.
On each atom the defining average gives
$\E[(X-M)\ind_B]=0$, hence $\E[(X-M)Z]=0$ for every such $Z$.
Taking $Z=M-\E X$ in the square expansion
$X-\E X=(X-M)+(M-\E X)$ proves the finite total-variance identity
\begin{equation}
 \Var(X)=\E\bigl[\Var(X\mid\mathcal G)\bigr]
       +\Var\bigl(\E[X\mid\mathcal G]\bigr).
\end{equation}
If some atoms have zero mass, assign arbitrary constant values on them.
The conditional identities then hold outside a measurable null event;
scalar expectation and variance identities remain exact and independent
of those choices, since null cells contribute zero. Thus uniqueness is
pointwise when every atom has positive mass, and otherwise only up to
these versions. None of this requires bounded real payoffs: the values
need only form a finite list in $K$.

\subsection{A finite betting interpretation}

A ticket on $A$ paying $\ind_A$ and priced at $P(A)$ has centered gain
$\ind_A-P(A)$. Any finite $K$-linear combination of such gains has
expectation zero. It cannot be strictly negative in every state: on the
finite partition generated by the tickets, a strictly negative value on
every nonempty cell gives a strictly negative expectation, since at least
one cell has positive mass. This is a finite no-sure-loss statement with
$K$-valued stakes and payoffs. It does not license infinitely many tickets
or an infinite series of stakes.

Conversely, write a quoted price as $q(A)$ and its centered ticket as
$T_A=\ind_A-q(A)$. If $q(\Omega)\ne1$, then $T_\Omega$ is a
nonzero constant. If disjoint events violate finite additivity, then
\[
 T_{A\cup B}-T_A-T_B=q(A)+q(B)-q(A\cup B)
\]
is a nonzero constant. Choosing its sign gives a strict sure loss.
If $q(A)>1$, buying $T_A$ loses in every state; if $q(A)<0$, selling it
does. Thus absence of a strict sure loss forces precisely normalization,
nonnegativity and finite additivity. It does not force regularity:
on a two-point full event space, the weights $(1,0)$ are coherent although
one nonempty event has price zero. Physical
interpretations of infinite stakes are a separate modeling issue; finite
coherence alone does not require such stakes to be practically available.

\section{Conditional shadows and finite scale hierarchies}
\label{fsp:sec:shadows}

\subsection{The unconditional shadow is not enough}

The function $P_0(A)=\st P(A)$ is a real finitely additive probability.
For $P_0(B)>0$, Proposition~\ref{fsp:prop:st} gives
\begin{equation}\label{fsp:eq:stcondition}
 \st P(A\mid B)=P_0(A\mid B).
\end{equation}
For $P_0(B)=0<P(B)$ the right-hand side is undefined, and no replacement
can be inferred from $P_0$ alone.

\begin{example}[Identical shadows, different rare-event answers]
\label{fsp:ex:shadowloss}
On three states let
\begin{equation}
 p=(1-2t,t,t),\qquad q=(1-4t,t,3t),\qquad B=\{2,3\},
 \quad t=\omega^{-1}.
\end{equation}
Both shadows are $(1,0,0)$, whereas
$P(\{2\}\mid B)=1/2$ and $Q(\{2\}\mid B)=1/4$.
Their total variation distance is only $2t$.
\end{example}

For a regular probability there is a richer real object:
\begin{equation}\label{fsp:eq:conditionalshadow}
 C(A\mid B)=\st\frac{P(A\cap B)}{P(B)},\qquad B\ne\varnothing.
\end{equation}
For each fixed $B$ it is a real finitely additive probability concentrated
on $B$. Whenever $D\ne\varnothing$ and $B\cap D\ne\varnothing$,
\begin{equation}\label{fsp:eq:popperproduct}
 C(A\cap B\mid D)=C(A\mid B\cap D)C(B\mid D).
\end{equation}
Indeed the identity holds before taking standard part, and both factors
are limited. This remains true if $C(B\mid D)=0$. Such conditional
systems are closely related to the Popper-function literature
\cite{halpern,brickhill}; no convention for conditioning on the empty
set is needed here. Importantly, the existence of all ratios in
\eqref{fsp:eq:conditionalshadow} is supplied by regularity, not by a
purported universal rule for $0/0$.

\subsection{Leading atoms determine every conditional real shadow}

\begin{theorem}[Conditional leading-scale skeleton]\label{fsp:thm:skeleton}
Let a finite regular Hahn probability have atom weights
\begin{equation}
 p_i=c_i t^{a_i}(1+\eta_i),\qquad c_i>0,\quad v(\eta_i)>0,
\end{equation}
where a zero error is allowed. For a nonempty $B$, put
$a(B)=\min_{i\in B}a_i$. Then
\begin{equation}\label{fsp:eq:skeleton}
 \st P(A\mid B)
 =\frac{\sum_{i\in A\cap B,\ a_i=a(B)}c_i}
        {\sum_{i\in B,\ a_i=a(B)}c_i}.
\end{equation}
Consequently all conditional shadows are determined by the ordered
partition of atoms into equal-$a_i$ groups and their positive leading
coefficients, up to a common positive rescaling within each group.
\end{theorem}
\begin{proof}
Divide numerator and denominator of \eqref{fsp:eq:condition} by $t^{a(B)}$.
The denominator has positive real standard part
$\sum_{i\in B,a_i=a(B)}c_i$. Terms with larger exponent are
infinitesimal, and terms with minimal exponent have standard part $c_i$.
Proposition~\ref{fsp:prop:st} now applies. Common rescaling of the coefficients
of the first group meeting $B$ cancels in the quotient.
\end{proof}

Equivalently, give each group the real probability proportional to its
$c_i$. To condition, select the first group assigning positive mass to
$B$ and condition that group's law. This is a particularly transparent
finite lexicographic representation of \emph{conditional shadows}.
It does not reconstruct the full surreal probability or all expected
utility comparisons. For example, $(1/2+t,1/2-t)$ has just the uniform
leading group but prefers the first-state payoff to the second-state
payoff. Higher coefficients can resolve ties even within one leading
group.

\subsection{A finite compression that retains real-act comparisons}

\begin{theorem}[Signed-row lexicographic compression]\label{fsp:thm:compression}
Let $p=(p_1,\ldots,p_n)$ be a finite Hahn probability. There are
$r\le n$ real row vectors $u_1,\ldots,u_r$ such that for every
$d\in\R^n$ the sign of $\sum_i p_i d_i$ is the sign of the first
nonzero entry in
\begin{equation}
 (u_1\cdot d,\ldots,u_r\cdot d),
\end{equation}
and the expectation is zero exactly when all these entries are zero.
The rows can be chosen from the actual coefficient rows of $p$ in
increasing exponent order.
\end{theorem}
\begin{proof}
The finite union $W$ of the supports of the $p_i$ is well ordered.
For $\gamma\in W$ let
$r_\gamma=([t^\gamma]p_1,\ldots,[t^\gamma]p_n)\in\R^n$.
Choose the first nonzero row. Repeatedly choose the earliest row outside
the span of the previously chosen rows, stopping when no such row
remains. Each choice increases dimension, so there are at most $n$ choices.
Every skipped row is in the span of chosen rows that precede it.
If $u_j$ is the first chosen row with nonzero pairing with $d$, all rows
before $u_j$ pair to zero, so its exponent is the first nonzero
coefficient of $\sum_i p_i d_i$. Its sign is therefore the sign of that
series. If all chosen rows annihilate $d$, all coefficient rows do.
\end{proof}

Applied to the difference of two real-valued acts, the theorem retains
all their preference comparisons, including infinitesimal tie breaking.
The rows need not be positive probabilities: normalization gives total
coefficient zero at every nonzero exponent, so signed rows are natural.
The theorem is not a claim that an arbitrary signed-row list is a
positive lexicographic probability system. Nor does it compress comparisons
of arbitrary surreal-valued acts: multiplying by such acts can introduce
new exponents and cancellations. It is a finite-dimensional statement
with an explicit domain of payoffs.

\section{Bayesian inference at every scale}
\label{fsp:sec:bayes}

\subsection{Exact posterior algebra and scale selection}

For finitely many hypotheses with positive prior probabilities $p_i$ and
likelihoods $\ell_i\ge0$, not all zero, define
\begin{equation}\label{fsp:eq:bayes}
 q_i=\frac{p_i\ell_i}{Z},\qquad Z=\sum_jp_j\ell_j.
\end{equation}
If the likelihood is an event probability then $\ell_i\le1$; if it is a
density or an unnormalized score, this upper bound is unnecessary.
Multiplying all likelihoods by one positive factor leaves the posterior
unchanged. Bayes' formula is exact for either interpretation when the
likelihood is part of a specified model.

\begin{theorem}[Minimum-plus Bayes with coefficient tie breaking]
\label{fsp:thm:scalebayes}
Suppose $p_i,\ell_i>0$ lie in a Hahn field. Write
$a_i=v(p_i)$, $b_i=v(\ell_i)$, $c_i=\lc(p_i)$,
$d_i=\lc(\ell_i)$, and
\begin{equation}
 m=\min_i(a_i+b_i),\qquad D=\sum_{i:a_i+b_i=m}c_id_i.
\end{equation}
Then
\begin{equation}\label{fsp:eq:scalebayes}
 v(q_i)=a_i+b_i-m,\qquad \lc(q_i)=c_id_i/D.
\end{equation}
In particular, the standard posterior is supported exactly on the
minimizers of $a_i+b_i$, with probabilities proportional to $c_id_i$.
Zero likelihoods may be omitted and yield zero posterior mass.
\end{theorem}
\begin{proof}
The product $p_i\ell_i$ has valuation $a_i+b_i$ and leading coefficient
$c_id_i>0$. Equation~\eqref{fsp:eq:positiveval} shows that $Z$ has leading
term $Dt^m$. Dividing each numerator by this leading term proves
\eqref{fsp:eq:scalebayes}. Nonminimal terms then have positive valuation,
while minimal terms have valuation zero and the displayed real residue.
\end{proof}

The valuation rule is minimum-plus arithmetic: multiply probabilities by
adding ranks, combine alternatives by taking a minimum, and normalize by
subtracting the winning rank. Leading real coefficients provide the
information needed when ranks tie. It is related to ranking models
\cite{spohn,hammond}, but the exponents here belong to an ordered abelian
group; their addition is not automatically ordinal addition.

\begin{example}[Rare evidence reverses the dominant hypothesis]
\label{fsp:ex:bayes}
Take the prior \eqref{fsp:eq:intro} and likelihoods $(t^3,t,1)$. The products
are proportional to $(t^3,t^2,t^2)$, so
\begin{equation}\label{fsp:eq:bayesexample}
 q=\left(\frac{t}{2+t},\frac1{2+t},\frac1{2+t}\right),
 \qquad \st q=(0,1/2,1/2).
\end{equation}
The evidence probability is
$(t^3+2t^2)/(1+t+t^2)>0$. Its standard part is zero. An analysis that
first discards infinitesimal prior and likelihood terms cannot perform
this update; it has erased exactly the information that the observation
selects.
\end{example}

Successive finite likelihood updates commute because their product
commutes. Under conditional independence, posterior odds equal prior odds
times the product of likelihood ratios. Neither the field algebra nor
this statement licenses an unspecified infinite product of likelihoods.

\subsection{How much precision does conditioning require?}

For finite laws $p,q$ define
\begin{equation}
 \delta=\TV(P,Q)=\max_A|P(A)-Q(A)|
 =\frac12\sum_i|p_i-q_i|.
\end{equation}
The equality follows by taking the subset of positive differences and
using $\sum_i(p_i-q_i)=0$; it is valid over any ordered field.

\begin{theorem}[Absolute and valuation-sensitive conditioning stability]
\label{fsp:thm:stability}
Suppose $P(B)\ge b>\delta$. Then $Q(B)>0$ and for every $A$,
\begin{equation}\label{fsp:eq:conditionalbound}
 |P(A\mid B)-Q(A\mid B)|\le\frac{2\delta}{b}.
\end{equation}
For Hahn probabilities, let $\beta=v(P(B))$ and $\lambda\ge0$.
If
\begin{equation}\label{fsp:eq:precision}
 v(p_i-q_i)>\beta+\lambda\quad\text{for every atom }i,
\end{equation}
then $Q(B)$ has the same leading term as $P(B)$ and
\begin{equation}
 v\bigl(P(A\mid B)-Q(A\mid B)\bigr)>\lambda.
\end{equation}
\end{theorem}
\begin{proof}
The first assumption gives $Q(B)\ge P(B)-\delta>0$.
Set $x=P(A\cap B)$, $y=Q(A\cap B)$, $d=P(B)$, and $e=Q(B)$.
Then
\begin{equation}
 \left|\frac{x}{d}-\frac{y}{e}\right|
 \le\frac{|x-y|}{d}+\frac{y|e-d|}{de}
 \le\frac{2\delta}{d},
\end{equation}
since $0\le y\le e$. This implies \eqref{fsp:eq:conditionalbound}.
For the valuation assertion write $e=d+\Delta_d$ and
$y=x+\Delta_x$. Finite sums in \eqref{fsp:eq:precision} give
$v(\Delta_d),v(\Delta_x)>\beta+\lambda$.
Thus $v(e)=v(d)=\beta$ and $e$ has the same leading coefficient.
Also $0\le x\le d$ implies $v(x)\ge\beta$ when $x\ne0$.
The numerator $d\Delta_x-x\Delta_d$ of the difference has valuation
strictly greater than $2\beta+\lambda$, while $de$ has valuation
$2\beta$. Division proves the assertion, including a zero numerator.
\end{proof}

Consequently an approximation suitable for unconditional probabilities
may be unusable after rare conditioning. To preserve a conditional answer
through scale $\lambda$, one needs input errors beyond the event scale
$\beta+\lambda$. Example~\ref{fsp:ex:shadowloss} demonstrates the qualitative
issue: an infinitesimal unconditional discrepancy can become a real
conditional discrepancy.

\subsection{Infinite prior odds are not ordinary dogmatism}

If $0<p<1$, a nonzero likelihood can always change its odds algebraically.
Nevertheless, an infinite prior logit cannot be brought to a finite logit
by an ordinary finite number of bounded-real log-likelihood increments.
If $L$ is positive infinite and $|b_j|\le C\in\R$ for $j\le n$ with
ordinary finite $n$, then $L+\sum_{j=1}^n b_j$ remains positive infinite.
Thus the posterior shadow remains one although the exact posterior changes.
A countervailing infinite-scale likelihood can overturn it. A statement
about learning from an infinite data stream requires a chosen probability
law on streams and a chosen limiting operation; it does not follow from
this finite observation.

\section{Surreal logits, log-scores, and softmax geometry}
\label{fsp:sec:logits}

\subsection{Binary probabilities and exact log-odds}

\begin{definition}
In an exponential/logarithmic workspace $E$, define
\begin{equation}\label{fsp:eq:logit}
 \sigma(z)=\frac{1}{1+\exp(-z)},\qquad
 \logit(p)=\log\frac{p}{1-p}\quad(0<p<1).
\end{equation}
\end{definition}

\begin{proposition}[Logit coordinates]
The maps $\sigma:E\to(0,1)_E$ and $\logit:(0,1)_E\to E$ are inverse,
strictly increasing bijections, and $\sigma(-z)=1-\sigma(z)$.
For binary hypotheses and positive likelihoods,
\begin{equation}\label{fsp:eq:logitbayes}
 \logit P(H\mid D)=\logit P(H)
 +\log\frac{P(D\mid H)}{P(D\mid H^c)}.
\end{equation}
\end{proposition}
\begin{proof}
For every $z$, $\sigma(z)/(1-\sigma(z))=\exp z$.
Conversely, substituting $p/(1-p)$ for $\exp z$ gives $\sigma(z)=p$.
Monotonicity and the complement identity follow from positivity and
monotonicity of $\exp$. Posterior odds are prior odds times the likelihood
ratio, and $\log(xy)=\log x+\log y$ proves \eqref{fsp:eq:logitbayes}.
\end{proof}

Every actual surreal $z$, including $\omega$ or $-\omega^2$, gives a
probability strictly between zero and one. The endpoints $p=0,1$ require
formal logit symbols $-\infty,+\infty$ outside $\No$, not the surreals
$-\omega,+\omega$. This distinction prevents an infinite confidence level
from being confused with a logically impossible alternative.

For positive infinite $L$ put $u=\exp(-L)$, which is a positive
infinitesimal. Then
\begin{equation}\label{fsp:eq:logittail}
 \sigma(-L)=\frac{u}{1+u}=u-u^2+u^3-\cdots,\qquad
 \sigma(L)=1-u+u^2-u^3+\cdots.
\end{equation}
The displayed series are strong normal-form substitutions in the ambient
surreals; a particular small workspace may contain the resulting rational
expression without being closed under arbitrary Hahn sums. For real $r$
and infinitesimal $\eta$, the canonical infinitesimal exponential expansion
gives
\begin{equation}
 \sigma(r+\eta)=\sigma(r)+\sigma(r)(1-\sigma(r))\eta+O(\eta^2).
\end{equation}
Here $O(\eta^2)$ means bounded by a real constant times $|\eta|^2$ in this
local expansion. For $t=\omega^{-1}$,
\begin{equation}\label{fsp:eq:logitt}
 \logit(t)=\log t-\log(1-t)
 =-\log\omega+t+\frac{t^2}{2}+\frac{t^3}{3}+\cdots.
\end{equation}
Thus a probability of order $\omega^{-1}$ corresponds to a logit near
$-\log\omega$, not near $-\omega$. A logit $-\omega$ specifies the much
smaller probability of order $\exp(-\omega)$.

\subsection{Logits do not inherit event additivity}

Logits are coordinates of probabilities, not a replacement additive
measure. If independent events have logits $a,b$, then
\begin{equation}\label{fsp:eq:andlogit}
 \logit P(A\cap B)=a+b-\log(1+e^a+e^b).
\end{equation}
Indeed, setting $u=e^a$, $v=e^b$ gives
$P(A\cap B)=uv/((1+u)(1+v))$, whose odds are $uv/(1+u+v)$.
For disjoint events, one adds $\sigma(a)+\sigma(b)$, when this is at most
one, and only then takes a logit if the result lies in $(0,1)$.
Adding $a+b$ is appropriate for independent \emph{evidence increments in
odds}, not for a disjoint union of events.

There is also a modeling distinction between a very small probability
and an unnormalized score. An isolated score becomes a probability only
relative to competing scores and a normalization rule. Surreal scores do
not remove that requirement.

\subsection{Softmax, gauge invariance, and stable normalization}

For an ordinary finite number of scores $z_1,\ldots,z_n\in E$, define
\begin{equation}\label{fsp:eq:softmax}
 \LSE(z)=\log\sum_{j=1}^n e^{z_j},\qquad
 p_i=e^{z_i-\LSE(z)}.
\end{equation}
All $p_i$ are positive and sum to one. Every strictly positive finite
law occurs, by taking $z_i=\log p_i$. Adding the same $c\in E$ to all
scores does not change $p$. For $n\ge2$, the one-against-the-rest logit is
\begin{equation}
 \logit(p_i)=z_i-\log\sum_{j\ne i}e^{z_j},
\end{equation}
not the raw score $z_i$ in general.

\begin{proposition}[Normalization and perturbation bounds]
\label{fsp:prop:softbound}
Let $M=\max_i z_i$. Then
\begin{equation}
 M\le\LSE(z)\le M+\log n.
\end{equation}
If $p$ and $q$ are the softmax laws of $z$ and $w$, and
$|z_i-w_i|\le\delta$ for every $i$, where $\delta\ge0$, then
\begin{equation}\label{fsp:eq:softbound}
 e^{-2\delta}p_i\le q_i\le e^{2\delta}p_i.
\end{equation}
\end{proposition}
\begin{proof}
Factor $e^M$ from the normalizing sum. The remaining positive sum lies
between $1$ and $n$. For the perturbation claim,
$e^{-\delta}e^{z_i}\le e^{w_i}\le e^\delta e^{z_i}$, and summing gives
the same bounds on the two normalizers. Dividing the bounds proves
\eqref{fsp:eq:softbound}.
\end{proof}

Infinitesimal score errors therefore produce infinitesimal relative errors
even in extremely small probabilities. This is a different guarantee from
an absolute error bound on probabilities, which can be destroyed by rare
conditioning. Subtracting $M$ before exponentiating is exact algebra and
avoids representing an unnecessarily large common exponential. It does
not by itself provide a finite symbolic representation for every surreal
score.

\subsection{Hierarchical scores and lexicographic concentration}

\begin{theorem}[Finite multiscale softmax limit]\label{fsp:thm:multisoft}
Let $\Omega_1,\ldots,\Omega_m$ be positive surreals with
$\Omega_m$ positive infinite and
$\Omega_{r+1}/\Omega_r$ infinitesimal for $r<m$. For finitely many states
suppose
\begin{equation}
 z_i=\sum_{r=1}^m\Omega_r a_{ir}+b_i+\eta_i,
 \quad a_{ir},b_i\in\R,\quad \eta_i\simeq0.
\end{equation}
Let $I_*$ be the states whose real tuple
$(a_{i1},\ldots,a_{im})$ is lexicographically maximal. Then
\begin{equation}\label{fsp:eq:multisoft}
 \st p_i=
 \begin{cases}
 e^{b_i}/\sum_{j\in I_*}e^{b_j},&i\in I_*,\\
 0,&i\notin I_*.
 \end{cases}
\end{equation}
\end{theorem}
\begin{proof}
Choose $j\in I_*$. If $i\notin I_*$, at the first differing coordinate
$r$ the coefficient of $\Omega_r$ in $z_i-z_j$ is a negative nonzero
real. All later contributions divided by $\Omega_r$ are infinitesimal,
and the bounded terms have that property too. Hence $z_i-z_j$ is
negative infinite and $e^{z_i-z_j}$ is infinitesimal. For $i,j\in I_*$,
all hierarchical terms cancel and
$\st e^{z_i-z_j}=e^{b_i-b_j}$. Divide the finite normalizing sum by
$e^{z_j}$ and take standard parts.
\end{proof}

This supplies an exact scalar model of ordered priorities, with ordinary
softmax resolving the final ties and smaller terms retaining further
information. For example, $\omega^2 a_i+\omega b_i+c_i$ first maximizes
$a_i$, then $b_i$, then uses $c_i$ as a finite log-score. Here $\omega^2$
is just the square of $\omega$, so the exponentiation warning in
Section~\ref{fsp:sec:scalars} is not implicated. The standard part describes
concentration; the exact law still gives positive mass to every state.

\section{Finite information theory and proper scoring}
\label{fsp:sec:information}

\subsection{Entropy and relative entropy}

All sums in this section have an ordinary finite number of terms and use
canonical logarithms. For probability vectors $p,q$ define
\begin{equation}\label{fsp:eq:entropy}
 H(p)=-\sum_{i:p_i>0}p_i\log p_i,\qquad
 D(p\Vert q)=\sum_{i:p_i>0}p_i\log\frac{p_i}{q_i},
\end{equation}
where the relative entropy is defined in $E$ if $q_i>0$ whenever $p_i>0$.
If that condition fails, it is useful to assign a formal value $+\infty$;
that value is not a surreal and is not used in field manipulations.
The convention $0\log0=0$ refers to entropy, not to an evaluation of
$\log0$.

\begin{theorem}[Finite Gibbs inequality and entropy bounds]
\label{fsp:thm:kl}
Under the support condition in \eqref{fsp:eq:entropy},
\begin{equation}
 D(p\Vert q)\ge0,
 \qquad D(p\Vert q)=0\ \Longleftrightarrow\ p=q.
\end{equation}
For $n$ states,
\begin{equation}\label{fsp:eq:entropybound}
 0\le H(p)\le\log n,
 \qquad \st H(p)=H(\st p).
\end{equation}
\end{theorem}
\begin{proof}
Let $S=\{i:p_i>0\}$. Apply $-\log x\ge1-x$ to $x=q_i/p_i$ and
multiply by $p_i>0$. Summing gives
\begin{equation}
 D(p\Vert q)\ge\sum_{i\in S}(p_i-q_i)
 =1-\sum_{i\in S}q_i\ge0.
\end{equation}
The difference between the first two expressions is a finite sum of
nonnegative terms, each zero only if $q_i=p_i$. Equality therefore
forces those equalities on $S$ and zero mass outside $S$, and the
converse is immediate. Each $-p_i\log p_i$ is nonnegative. Applying
Gibbs inequality to the uniform law gives
$D(p\Vert(1/n,\ldots,1/n))=\log n-H(p)\ge0$.

For the standard-part identity, if $\st p_i=r_i>0$, then
$\log p_i\simeq\log r_i$ by the infinitesimal logarithm expansion.
If $p_i>0$ is infinitesimal, the real-exponential inequality
$\log x\le\sqrt{x}$ for real $x\ge4$ holds for surreal $x\ge4$ too.
Consequently
$0\le p_i\log(1/p_i)\le\sqrt{p_i}\simeq0$.
Taking standard parts of the finite sum proves the formula, including
zero coordinates.
\end{proof}

An entropy on a fixed finite alphabet is always limited. It nevertheless
retains infinitesimal information discarded by its standard part. For a
rare Bernoulli success of probability $t$,
\begin{equation}\label{fsp:eq:rareentropy}
 H(t,1-t)=t\log(1/t)+t-\frac{t^2}{2}+O(t^3).
\end{equation}
The log factor is essential: entropy need not remain in the original
pure Laurent field $\R((t))$. It belongs to the exponential/logarithmic
workspace, where the displayed products are meaningful. This is one
reason not to assume that a single Hahn field is automatically closed
under all information-theoretic operations.

\subsection{Conditional information and data processing}

For a finite joint law write $p_{ij}=p^Y_jp^{X\mid j}_i$ on every cell
with $p^Y_j>0$. Expanding $\log(p^Y_jp^{X\mid j}_i)$ and summing yields
\begin{equation}
 H(X,Y)=H(Y)+\sum_jp^Y_jH(X\mid Y=j),
\end{equation}
where zero-mass cells contribute zero. The same decomposition for two
joint laws satisfying the relative-entropy support condition gives
\begin{equation}\label{fsp:eq:klchain}
 D(p_{XY}\Vert q_{XY})
 =D(p_Y\Vert q_Y)
 +\sum_{j:p^Y_j>0}p^Y_j
       D(p_{X\mid j}\Vert q_{X\mid j}).
\end{equation}
To verify it, replace every logarithm of a joint ratio by the logarithm
of the marginal ratio plus that of the conditional ratio. All terms
are finite sums, so regrouping requires no additional hypothesis.

\begin{theorem}[Finite data processing]\label{fsp:thm:dataprocessing}
Let $T_{ij}\ge0$ be a finite stochastic matrix, $\sum_jT_{ij}=1$.
For $p,q$ satisfying the support condition,
\begin{equation}
 D(pT\Vert qT)\le D(p\Vert q).
\end{equation}
Equality holds exactly when, for each output $j$ with $(pT)_j>0$,
the two input conditional laws given $j$ coincide.
\end{theorem}
\begin{proof}
Form the joint laws $P_{ij}=p_iT_{ij}$ and $Q_{ij}=q_iT_{ij}$.
On each positive $P_{ij}$ their ratio is $p_i/q_i$; summing over $j$
gives $D(P\Vert Q)=D(p\Vert q)$. Their output marginals are $pT,qT$.
Equation~\eqref{fsp:eq:klchain} and Theorem~\ref{fsp:thm:kl} give the inequality.
The difference is a finite sum with positive weights, so it vanishes
exactly when every indicated conditional divergence vanishes.
\end{proof}

Taking $q_{ij}=p^X_ip^Y_j$ gives the nonnegative mutual information
$I(X;Y)=D(p_{XY}\Vert p_Xp_Y)$ and the usual finite identities
$I(X;Y)=H(X)+H(Y)-H(X,Y)$ and
$I(X;Y)=H(X)-H(X\mid Y)$. Independence is equivalent to zero mutual
information. A process-level entropy rate would involve a limit in the
horizon and is not supplied by these finite formulas.

\subsection{Identical real shadows can hide infinite divergence}

\begin{example}[An information scale invisible to standard probability]
\label{fsp:ex:kl}
Let $t=\omega^{-1}$ and $u=\exp(-1/t^2)$. Consider
\begin{equation}
 p=(t,1-t),\qquad q=(u,1-u).
\end{equation}
Both standard parts equal $(0,1)$, but
\begin{align}
 D(p\Vert q)
 &=\frac1t+t\log t
   +(1-t)\bigl(\log(1-t)-\log(1-u)\bigr),\label{fsp:eq:hugekl}\\
 \st\bigl(tD(p\Vert q)\bigr)&=1.
\end{align}
Thus their relative entropy is positive infinite. In fact the final
term in \eqref{fsp:eq:hugekl} is infinitesimal, and $t\log t$ is
infinitesimal, while the first term is $1/t$.
\end{example}

There is no conflict with \eqref{fsp:eq:entropybound}. Entropy uses each
law's own probabilities and is bounded on a fixed alphabet. Relative
entropy can compare a rare event against a vastly smaller prediction.
Taking standard parts before forming the logarithmic ratio loses this
penalty. Standard part therefore commutes with finite entropy but does
not in general commute with relative entropy, even when the latter
has a well-defined surreal value.

\subsection{Strictly proper scores survive}

For an outcome $i$, logarithmic loss is $-\log q_i$. If all predicted
probabilities are positive, its expected excess over truthful prediction
is exactly
\begin{equation}
 \E_p[-\log q_i]-H(p)=D(p\Vert q).
\end{equation}
On the strictly positive simplex it is therefore uniquely minimized by
$q=p$, by Theorem~\ref{fsp:thm:kl}. The Brier loss
$S(q,i)=\sum_j(q_j-\ind_{j=i})^2$ satisfies the entirely algebraic identity
\begin{equation}
 \E_p S(q,i)-\E_p S(p,i)=\sum_j(q_j-p_j)^2.
\end{equation}
It is strictly proper on the whole finite simplex over any ordered field.
These scores distinguish infinitesimal prediction errors, though their
relative penalties can occur at different scales.

For a binary soft target $0<p<1$ and logit $z$, the expected logarithmic
loss is
\begin{equation}\label{fsp:eq:logisticloss}
 \Phi(z;p)=\log(1+e^z)-pz.
\end{equation}
Since $q=\sigma(z)$ parametrizes the strictly positive binary simplex,
its unique minimizer is $z=\logit(p)$. This conclusion uses an explicit
identity and positivity, not an unproved extreme-value theorem on a
surreal parameter space.

\section{Gibbs laws, decision theory, and finite learning models}
\label{fsp:sec:gibbs}

\subsection{A finite variational principle without compactness}

\begin{theorem}[Surreal finite Gibbs principle]\label{fsp:thm:gibbs}
Let $E_1,\ldots,E_n\in E$ be energies and let $\tau>0$ belong to $E$.
Set
\begin{equation}
 Z=\sum_i e^{-E_i/\tau},\qquad q_i=Z^{-1}e^{-E_i/\tau}.
\end{equation}
For every probability vector $p$,
\begin{equation}\label{fsp:eq:gibbs}
 \sum_i p_iE_i-\tau H(p)
 =-\tau\log Z+\tau D(p\Vert q).
\end{equation}
Hence $q$ is the unique minimizer of the left-hand side on the finite
simplex, even when some energies or $1/\tau$ are infinite.
\end{theorem}
\begin{proof}
Use $\log q_i=-E_i/\tau-\log Z$ in the definition of relative entropy.
Multiplication by $\tau$ and rearrangement give \eqref{fsp:eq:gibbs}.
The last term is nonnegative and vanishes exactly at $p=q$.
\end{proof}

No compactness or Dedekind completeness of the simplex was invoked.
The minimizer was written down, and the identity proved optimality.
Infinitesimal temperature is meaningful for a finite system: among real
energies with a positive real gap, excited states have infinitesimal
probability. Several ordered energy scales can be treated by
Theorem~\ref{fsp:thm:multisoft}. This does not construct a thermodynamic
limit, an infinite-volume Gibbs measure, or a partition function over
an arbitrary infinite state space.

\subsection{Expected utility and the relevance of rare events}

For finitely many acts and states, expected utility is an ordered-field
calculation. A maximum among finitely many expected utilities exists.
No Archimedean continuity of preferences is forced: a positive
infinitesimal preference can be weaker than every positive real
preference and still be strict. The compression in
Theorem~\ref{fsp:thm:compression} explains how real-valued acts can detect
successive signed coefficient layers.

There is an additional reason not to delete rare events when payoffs
are allowed to be large. An event of probability $t$ with loss $1/t$
has expected loss exactly one. Replacing its probability by zero before
multiplication gives the wrong answer. By contrast, for finitely many
\emph{limited} payoffs $X_i$, standard part obeys
\begin{equation}
 \st\E_pX=\sum_i(\st p_i)(\st X_i),
\end{equation}
by its ring property. The qualification on the payoffs is indispensable.
A normative application must decide which utilities and trades it permits;
surreal arithmetic does not itself settle that decision.

\subsection{Infinitesimal smoothing}

Suppose a finite alphabet has ordinary counts $N_i\ge0$ with
$N=\sum_iN_i>0$, and choose positive real constants $a_i$.
The law
\begin{equation}\label{fsp:eq:smoothing}
 p_i=\frac{N_i+t a_i}{N+t\sum_j a_j}
\end{equation}
is regular. Seen categories have standard probabilities $N_i/N$;
unseen categories have positive infinitesimal probabilities with leading
coefficient $a_i/N$ at scale $t$. Subsequent finite Bayes updates are
well defined on all categories. The constants $a_i$ and the scale $t$
are modeling choices, not consequences of the data. Distinct choices
can give the same empirical shadow and different answers after rare
conditioning.

For $N=0$, formula \eqref{fsp:eq:smoothing} instead reduces to the ordinary
prior $a_i/\sum_j a_j$; this case should be handled explicitly rather
than described as an infinitesimal perturbation of absent empirical
frequencies.

\subsection{Surreals do not automatically cure separation}

\begin{proposition}[No attained logistic optimum under strict separation]
\label{fsp:prop:separation}
For a nonempty finite collection of feature vectors $x_i\in E^d$ and labels
$y_i\in\{-1,1\}$, define
\begin{equation}
 L(\theta)=\sum_i\log(1+\exp(-y_i x_i\cdot\theta)).
\end{equation}
If there exists $v\in E^d$ with $y_i x_i\cdot v>0$ for every $i$,
then $L$ has no minimizer in $E^d$.
\end{proposition}
\begin{proof}
For every $s>0$, each signed margin strictly increases from $\theta$
to $\theta+sv$. Hence every exponential in the sum strictly decreases,
and so does every logarithmic loss. Their finite sum strictly decreases.
This improves every proposed parameter, including parameters with
infinite coordinates.
\end{proof}

An infinite parameter is not a terminal point at infinity: it can still
be increased. Infinitesimal positive losses are not zero. Soft targets
with independent logits have the explicit optimum described in
\eqref{fsp:eq:logisticloss}; that fact does not establish existence for every
constrained regression or regularization problem. Such problems need
their own hypotheses and proofs.

\section{Finite stochastic processes and the time-scale issue}
\label{fsp:sec:processes}

\subsection{Finite chains and finite paths}\label{fsp:sub:paths}

Let $T$ be a finite stochastic matrix over $K$ and let $\pi$ be an initial
law. For each ordinary finite horizon $N$ the path weights
\begin{equation}
 P(i_0,\ldots,i_N)=\pi_{i_0}\prod_{k=1}^N T_{i_{k-1},i_k}
\end{equation}
are nonnegative and normalized. Summing the last coordinate proves
consistency when the horizon is shortened. Standard part commutes with
all these finite products and sums, so the path shadow is the real
chain with initial law $\st\pi$ and transition matrix $\st T$.
This consistency alone is not a theorem about a probability measure on
infinite paths; Section~\ref{fsp:sec:iid} supplies concrete obstructions to
one proposed countable-additivity rule.

The repository's analysis of finite Hahn Markov generators and their
resolvents concerns another important part of multiscale stochastic
structure \cite{repomarkov}. A stochastic matrix at finitely many steps,
a continuous-time semigroup, and a path measure are different objects.
This article only asserts a path construction when the relevant
finite or infinite probability framework has been specified.

\subsection{Finite martingales and stopping}\label{fsp:sub:stopping}

On a finite probability space with a finite filtration
$\F_0\subseteq\cdots\subseteq\F_N$, a martingale satisfies
$\E[M_k\mid\F_{k-1}]=M_{k-1}$. The conditional expectation is the one
already constructed by finite partitions.

\begin{proposition}[Bounded finite-horizon optional stopping]
If $T$ is a stopping time taking values in $\{0,\ldots,N\}$, then
$\E M_T=\E M_0$.
\end{proposition}
\begin{proof}
The pathwise identity
\begin{equation}
 M_T=M_0+\sum_{k=1}^N\ind_{\{T\ge k\}}(M_k-M_{k-1})
\end{equation}
is a finite telescoping sum. The indicator is
$\F_{k-1}$-measurable. By the tower and pull-out properties, each
summand has expectation zero. Taking expectations proves the claim.
\end{proof}

No boundedness by a real constant of $M_k$ is needed on this finite
space. Conversely, this proof does not justify an unbounded stopping
time or a limit of stopped martingales. Those assertions require
additional integration and convergence results.

\subsection{Rare transitions and sample size}

Consider a state that changes irreversibly with probability $t>0$
infinitesimal at each step. At a fixed ordinary horizon $n$ its hitting
probability is
\begin{equation}\label{fsp:eq:hitting}
 1-(1-t)^n\le nt.
\end{equation}
The inequality follows by the finite union bound, or by induction on
$n$. Thus every ordinary finite-horizon shadow sees no transition,
although the exact finite-horizon probability is positive.
Writing ``take $n=1/t$'' is not a valid use of this formula unless a
meaning for a nonstandard or surreal-indexed number of trials has first
been supplied. A scalar in $\No$ is not automatically a finite set
cardinality, a path length, or an induction index.

For $n$ independent Bernoulli variables at an ordinary finite horizon,
finite variance algebra gives
\begin{equation}
 \E\overline X_n=p,\qquad
 \Var(\overline X_n)=\frac{p(1-p)}{n},\qquad
 P(|\overline X_n-p|\ge a)\le\frac{p(1-p)}{na^2}.
\end{equation}
These inequalities hold for surreal $a>0$. Their interpretation depends
on scale. For fixed positive real $a$ and a real shadow, they have the
usual ordinary limiting implications in a compatible classical model.
For an infinitesimal resolution $a$, increasing an ordinary integer
$n$ need not make the right-hand side small. There is no unconditional
law of large numbers in the full fine surreal topology hidden in this
finite calculation.

\section{Why infinite addition is an additional axiom}
\label{fsp:sec:infinity}

\subsection{The fine topology is too strong for ordinary probabilistic limits}

\begin{proposition}[Set-indexed fine convergence is eventually constant]
\label{fsp:prop:fine}
If a net $(x_i)_{i\in I}$ with set-sized index set converges to $x$ in
the full ordered-field topology of $\No$, then $x_i=x$ eventually.
\end{proposition}
\begin{proof}
The nonzero distances form a set
$D=\{|x_i-x|:x_i\ne x\}$ of positive surreals. If $D$ is empty the
claim is immediate. Otherwise the surreal cut property gives a positive
$\delta$ smaller than every member of $D$, for example a number between
$\{0\}$ and $D$. Convergence makes $|x_i-x|<\delta$ eventually. By
construction this forces $x_i=x$.
\end{proof}

The statement concerns the full fine topology and set-sized indices.
It does not concern an intrinsic topology on a fixed subfield: the
$\delta$ chosen in the proof may lie outside that subfield. Nevertheless,
even in any ordered subfield containing a positive infinitesimal $t$,
the ordinary real partial sums of $\sum_{n\ge1}2^{-n}$ do not converge
to one in the field's order topology. Their real tails $2^{-N}$ exceed
$t$ for every ordinary $N$. Coefficientwise real convergence and intrinsic
valuation convergence are different again. These distinctions agree with
the repository's foundations discussion \cite{repofound}, and
Proposition~\ref{fsp:prop:fine} is clause~(2) of a theorem of that report
(Section~\ref{fsp:sec:neighbours}).

Thus ordinary countable additivity cannot be imported by writing an
unqualified order-topological limit of finite partial sums. A successful
infinite theory must choose a different addition convention, a different
topology, or a standard-part procedure.

\subsection{What strong Hahn addition does and does not allow}

A strong Hahn probability on a sigma-algebra requires, for every disjoint
sequence $(A_n)$, that $(P(A_n))$ be strongly Hahn summable and that
\begin{equation}
 P\Bigl(\bigcup_n A_n\Bigr)=\hsum_nP(A_n).
\end{equation}
The ordinary series $\sum_n2^{-n}$ is not strong: every summand contributes
to exponent zero. An infinite collection of equal positive infinitesimal
singleton masses is not strong either, since the leading monomial repeats
infinitely often. These are failures of a specified support rule, not
contradictions in the ordered field or in finite additivity.

\begin{example}[A regular strong Hahn probability on a countably infinite space]
\label{fsp:ex:geometric}
On $\N_{\ge1}$ take $t=\omega^{-1}$ and
\begin{equation}
 w_n=(1-t)t^{n-1},\qquad
 P(A)=\hsum_{n\in A}w_n.
\end{equation}
Every exponent occurs in at most two weights. All subfamilies are
strongly summable, and the sum of all weights is one by coefficientwise
telescoping. For nonempty $A$, its least element $m$ supplies the leading
term $t^{m-1}$ with coefficient one, so $P(A)>0$. Disjoint regrouping
is still strong, proving countable additivity in this sense. The standard
part is the point mass at $1$.
\end{example}

Strong probabilities on countably separated spaces are much more rigid
than arbitrary finite-additive or coefficientwise laws. The repository's
classification gives finitely atomic coefficient measures and a finitely
supported standard part in that class \cite{repomeasures}
(Section~\ref{fsp:sec:neighbours}, item~(b)). We use the
explicit example and the elementary support obstructions above without
requiring the full classification proof. The real Lebesgue law, embedded
as exponent zero, is not strong because disjoint positive-measure intervals
make that exponent recur infinitely often.

\subsection{Coefficientwise countable additivity}

\begin{definition}\label{fsp:def:coefficientwise}
A coefficientwise Hahn measure on $(\Omega,\A)$ is specified by a
well-ordered set $W\subseteq\Gamma$ and ordinary finite signed measures
$\mu_\gamma$ on $\A$, via
\begin{equation}\label{fsp:eq:coefficientmeasure}
 P(A)=\sum_{\gamma\in W}\mu_\gamma(A)t^\gamma.
\end{equation}
It is a probability when these Hahn values are nonnegative on every
$A$ and $P(\Omega)=1$. Countable additivity means that each coefficient
uses the ordinary absolutely convergent scalar sum on a disjoint sequence.
\end{definition}

The definition assumes a sigma-algebra and ordinary countably additive
finite signed measures. It is strictly more permissive than the strong
rule in familiar separated spaces: real probability measures, diffuse or
atomic, are included as single exponent-zero coefficients. Higher
coefficient measures may be signed even when the reassembled Hahn
probability is positive. Positivity of every cylinder in a product or
of every finite test does not automatically imply positivity on all
measurable sets; the repository treats this issue explicitly
\cite{repomeasures}. In the constructions below, positivity is proved
directly rather than inferred from finite tests alone.

\section{Diffuse hierarchical probabilities and integration}
\label{fsp:sec:hierarchy}

\subsection{A positive construction with arbitrarily many ordered layers}

\begin{theorem}[Normalized hierarchy of real probability laws]
\label{fsp:thm:hierarchy}
Let $(\mu_\alpha)_{\alpha<\kappa}$ be a nonempty set-sized well-ordered
family of ordinary real probability measures on one measurable space.
Choose strictly increasing exponents $\gamma_\alpha\in\Gamma$ with
$\gamma_0=0$. Define
\begin{equation}\label{fsp:eq:hierarchy}
 Z=\hsum_{\alpha<\kappa}t^{\gamma_\alpha},\qquad
 P(A)=Z^{-1}\hsum_{\alpha<\kappa}t^{\gamma_\alpha}\mu_\alpha(A).
\end{equation}
Then $P$ is a coefficientwise Hahn probability. An event $B$ has positive
mass exactly when some $\mu_\alpha(B)>0$. If $\alpha(B)$ is the first
such index, then
\begin{equation}\label{fsp:eq:hierconditional}
 \st P(A\mid B)
 =\frac{\mu_{\alpha(B)}(A\cap B)}{\mu_{\alpha(B)}(B)}.
\end{equation}
In particular $\st P=\mu_0$.
\end{theorem}
\begin{proof}
The monomial family has a well-ordered support and one contribution at
each exponent, so both displayed sums exist. The denominator has leading
term one and is positive. A nonzero numerator has a first positive
coefficient, so it is positive. The numerator for $\Omega$ is $Z$,
giving normalization; finite additivity follows coefficientwise.

To verify the full coefficientwise definition, expand $Z^{-1}$ as a Hahn
series. Its support is well ordered by the positive-support Neumann
lemma, and multiplying by the common support
$\{\gamma_\alpha\}$ has only finitely many contributions at each exponent.
Thus each coefficient of $P$ is a finite real linear combination of the
$\mu_\alpha$, hence an ordinary finite signed measure. The resulting
support is common to all events. This proves coefficientwise countable
additivity. Finally the factor $Z^{-1}$ cancels in conditional ratios;
dividing by $t^{\gamma_{\alpha(B)}}$ and taking residues proves
\eqref{fsp:eq:hierconditional}.
\end{proof}

This construction is not restricted to finite or countable hierarchies.
Its existence uses set-sized well-ordered support, not a real-convergent
sum of the coefficients of $Z$. In particular, one must not evaluate
$t=1$ in \eqref{fsp:eq:hierarchy} and regard the generally divergent result
as an ordinary mixture.

\begin{example}[Diffuse probability with exceptional point events]
\label{fsp:ex:diffuse}
On the Borel subsets of $[0,1]$ let $\lambda$ be Lebesgue probability and
put
\begin{equation}
 P(A)=\frac{\lambda(A)+t\delta_0(A)+t^2\delta_1(A)}{1+t+t^2}.
\end{equation}
Its standard part is exactly Lebesgue probability, but both endpoint
singletons have positive infinitesimal mass. Conditioning on
$\{0,1\}$ gives masses $1/(1+t)$ and $t/(1+t)$, with conditional shadow
concentrated at zero. Other singletons retain exact probability zero.
Thus this is not an all-subsets regular probability, and it does not
solve conditioning on every exact-null event.
\end{example}

Higher coefficients of this normalized law are signed because division
by $1+t+t^2$ introduces corrections. Those signed coefficients are
compatible with positivity of every event, which was proved before
expanding the reciprocal. This is a useful construction pattern: prove
positivity of the normalized expression rather than demand positivity
of every coefficient.

\subsection{A controlled expectation for bounded real observables}

For the hierarchy in Theorem~\ref{fsp:thm:hierarchy} and a bounded measurable
\emph{real-valued} function $f$, define
\begin{equation}\label{fsp:eq:hierintegral}
 \E_P f=Z^{-1}\hsum_{\alpha<\kappa}t^{\gamma_\alpha}
                  \int f\,d\mu_\alpha.
\end{equation}
This is a positive real-linear functional, agrees with event probabilities
on indicators, and obeys $|\E_Pf|\le\norm{f}_\infty$. Indeed, if $f\ge0$
all integrals in the numerator are nonnegative; apply this observation
to $\norm{f}_\infty\pm f$ for the bound. It also gives
$\st\E_Pf=\int f\,d\mu_0$.

If $f_n$ are uniformly bounded real functions converging pointwise
outside a null set for every component $\mu_\alpha$ to $f$, ordinary
dominated convergence gives convergence of every component integral.
Every coefficient of the normalized series is a finite linear
combination of component integrals. Thus $\E_Pf_n$ converges
\emph{coefficientwise in the real coefficient topology} to $\E_Pf$.
This is a precise dominated-convergence statement, but it is not fine
or ordered-field convergence. For example, a nonconstant ordinary real
sequence of expectations need not converge in the order topology of a
field with infinitesimals.

For finite hierarchies on two spaces, form products of their component
measures. The weights are products of monomials, and the normalizer is
the product of normalizers. Ordinary Fubini for a bounded real function
on each component product then proves the corresponding coefficientwise
Fubini identity. More explicitly, for two hierarchies with exponents
$\gamma_\alpha,\delta_\beta$, the iterated integral is defined through
component pairing as
\begin{equation}
 \frac{1}{Z_PZ_Q}\hsum_{\alpha,\beta}
 t^{\gamma_\alpha+\delta_\beta}
 \int\!\int f(x,y)\,d\nu_\beta(y)\,d\mu_\alpha(x).
\end{equation}
One does not apply the real-observable definition unchanged to an arbitrary
Hahn-valued inner integral. The same argument applies to longer hierarchies when
the two support families satisfy the Hahn product support condition,
which guarantees finitely many contributions to each coefficient.
This is a support-controlled construction, not a Fubini theorem for all
possible surreal-valued functions.

\subsection{A genuine obstruction for arbitrary surreal-valued observables}

A pointwise Hahn-valued function can have one well-ordered support and
measurable real coefficients while naive coefficientwise integration
fails positivity. The distinction between real-valued observables and
arbitrary surreal-valued observables is therefore substantive.

\begin{theorem}[An uncountable-support positivity failure]
\label{fsp:thm:integrationfailure}
There are an ordinary countably additive probability $\mu$ and a
coefficientwise measurable Hahn-valued function $X$ with a common
well-ordered support such that $0<X<1$ everywhere, all coefficient
functions are bounded and integrable, but integrating each coefficient
separately gives a negative Hahn value.
\end{theorem}
\begin{proof}
Let the sample space be the first uncountable ordinal $\omega_1$, with
the countable--cocountable sigma-algebra. Give countable sets measure
zero and cocountable sets measure one. This is countably additive:
a disjoint sequence either has no cocountable member, in which case
its union is countable, or has exactly one, with all remaining members
countable.

Choose a set-sized ordered subgroup $\Gamma\subset\No$ containing all
ordinals $\xi\le\omega_1$, and use the associated Hahn field. Set
\begin{equation}\label{fsp:eq:badintegral}
 X(\xi)=t^\xi-t^{\omega_1}\quad(\xi<\omega_1).
\end{equation}
Because $\xi<\omega_1$, the first term is larger than the second;
and $t^\xi\le1$. Thus $0<X(\xi)<1$. The common exponent set
$\{\xi:\xi<\omega_1\}\cup\{\omega_1\}$ is well ordered.
At exponent $\alpha<\omega_1$ the coefficient function is the indicator
of $\{\alpha\}$; at exponent $\omega_1$ it is the constant $-1$.
All are measurable and bounded. Every singleton coefficient integrates
to zero, while the last coefficient integrates to $-1$. The formal
coefficientwise integral is consequently $-t^{\omega_1}<0$.
\end{proof}

The example asserts coefficientwise measurability, not measurability
for an unspecified topology on the Hahn field. It exhibits precisely
why that weak measurable-function class cannot support a positive
coefficientwise integration operation. It does not contradict ordinary
Lebesgue integration, the finite-range construction, or
\eqref{fsp:eq:hierintegral} for real-valued functions.

There is a useful boundary to this failure. If a pointwise nonnegative
Hahn-valued function has a \emph{countable} common well-ordered support
and each coefficient is integrable against an ordinary positive finite
measure, its coefficientwise integral is nonnegative. To see this,
proceed along the support. Provided all earlier coefficient integrals
are zero, induction makes all earlier coefficient functions zero
almost everywhere. Their predecessors are countable, so their common
zero set still has full measure. Pointwise Hahn positivity forces the
current coefficient to be nonnegative there. Its integral is
nonnegative, and if zero the coefficient also vanishes almost everywhere.
The first nonzero integral is therefore positive, or all integrals
are zero. The failure in \eqref{fsp:eq:badintegral} exploits an uncountable
union of null sets; that union need not be null.

\subsection{Posterior kernels in a finite diffuse hierarchy}
\label{fsp:subsec:kernels}

A conditional distribution for a continuous observation needs more than
a ratio of event masses at a singleton. Here is a construction with its
integration convention made explicit. Let
$\mu_0,\ldots,\mu_m$ be ordinary joint probability laws of $(X,Y)$ on
standard Borel spaces. Suppose their ordinary conditional kernels
$K_j(y,A)$ and marginal measures $\mu_j^Y$ have been selected, so that
\begin{equation}
 \mu_j(X\in A,Y\in B)=\int_B K_j(y,A)\,d\mu_j^Y(y).
\end{equation}
The usual real disintegration theorem supplies this input in this
setting \cite{kallenberg}; it is not a theorem about a new surreal
sigma-algebra.

Take $\rho=(m+1)^{-1}\sum_j\mu_j^Y$ and let
$h_j=d\mu_j^Y/d\rho$. Choose versions which are finite and nonnegative
on a common full-$\rho$ set. There $\sum_jh_j=m+1$, so at least one
$h_j$ is positive. For $0=\gamma_0<\cdots<\gamma_m$, put
\begin{align}
 D(y)&=\sum_{j=0}^m t^{\gamma_j}h_j(y),\label{fsp:eq:kerneldensity}\\
 P_y(A)&=\frac{\sum_{j=0}^m t^{\gamma_j}h_j(y)K_j(y,A)}{D(y)}.
 \label{fsp:eq:posterior}
\end{align}
For each such $y$ this is a positive normalized coefficientwise Hahn probability in
$A$. Its standard part selects the first $j$ with $h_j(y)>0$ and uses
$K_j(y,\cdot)$. One can assign any fixed probability kernel on the
remaining $\rho$-null set.

The precise reconstruction identity is
\begin{equation}\label{fsp:eq:kernelreconstruction}
 P(X\in A,Y\in B)
 =Z^{-1}\int_B D(y)P_y(A)\,d\rho(y)
 =Z^{-1}\sum_{j=0}^m t^{\gamma_j}
             \int_B h_j(y)K_j(y,A)\,d\rho(y),
\end{equation}
where $Z=\sum_jt^{\gamma_j}$. In the middle expression the product
$D(y)P_y(A)$ is formed \emph{before} coefficientwise integration;
it cancels the denominator, giving exactly the finite sum on the right.
Those coefficient functions are integrable, since $0\le K_j\le1$ and
$h_j$ is a probability density. Coefficients of the quotient
$P_y(A)$ alone need not have the integrability needed for an arbitrary
termwise integration operation.

Equation~\eqref{fsp:eq:kernelreconstruction} is a rigorous posterior
construction for this model, not a general surreal Radon--Nikodym or
disintegration theorem. Its version choices on sets null for every
component remain real choices. Replacing scalars does not determine a
unique conditional law on observations which the entire model assigns
exactly zero mass.

\section{Two coherent coin models with no coefficientwise path law}
\label{fsp:sec:iid}

For each ordinary finite $n$, independent coin weights are legitimate
surreal probabilities. The following proofs show why this is not enough
for the coefficientwise extension in Definition~\ref{fsp:def:coefficientwise}.
They concern $t=\omega^{-1}$ and the ordinary countable path space
$\{0,1\}^{\N}$ with its Borel sigma-algebra. The conclusions are special
cases of the repository's broader product analysis \cite{repomeasures};
Section~\ref{fsp:sec:neighbours}, item~(c), names the corollary of that report
which contains both and the mechanisms the proofs share.

\begin{theorem}[A repeated infinitesimal success probability]
\label{fsp:thm:rarecoin}
The consistent finite-dimensional laws of independent Bernoulli$(t)$
coordinates do not extend to a coefficientwise Hahn probability with
ordinary finite signed coefficient measures.
\end{theorem}
\begin{proof}
Let $B_n$ be the cylinder event whose first success occurs at time $n$.
These events are pairwise disjoint, and the prescribed finite-dimensional
laws require
\begin{equation}\label{fsp:eq:firstsuccess}
 P(B_n)=t(1-t)^{n-1}.
\end{equation}
The coefficient of $t$ is one for every $n$. If a coefficientwise law
existed, its exponent-one coefficient would be an ordinary finite
signed measure $\mu_1$ with $\mu_1(B_n)=1$ for all $n$. For any positive
integer $N$, its total variation would be at least
$\sum_{n=1}^N|\mu_1(B_n)|=N$, contradicting finiteness.
\end{proof}

The obstruction does not say that any finite coin experiment is
inconsistent. Nor does it say that a regular finitely additive path law
is impossible: Section~\ref{fsp:sec:extension} constructs such laws after
field enlargement. It says that one specific coefficientwise rule cannot
accommodate these finite-dimensional distributions.

\begin{theorem}[A repeated infinitesimal perturbation of a fair coin]
\label{fsp:thm:faircoin}
The independent Bernoulli$(1/2+t)$ finite-dimensional distributions also
have no coefficientwise extension with finite signed coefficient measures.
\end{theorem}
\begin{proof}
A length-$n$ cylinder containing $k$ successes has probability
$(1/2+t)^k(1/2-t)^{n-k}$. Its coefficient at $t$ is
\begin{equation}\label{fsp:eq:coinrow}
 2(2k-n)2^{-n}.
\end{equation}
The total variation of the proposed exponent-one coefficient measure,
evaluated on the length-$n$ partition, must therefore be at least
$2\E|S_n|$, where $S_n$ is a sum of $n$ ordinary independent symmetric
signs. Ordinary finite moment expansion gives
$\E S_n^2=n$ and $\E S_n^4=3n^2-2n\le3n^2$.
H\"older applied to
$|S_n|^2=|S_n|^{2/3}|S_n|^{4/3}$ gives
\begin{equation}
 n\le(\E|S_n|)^{2/3}(\E S_n^4)^{1/3},
 \qquad \E|S_n|\ge\sqrt{n/3}.
\end{equation}
The variation lower bound tends to infinity through ordinary real
numbers. A finite signed coefficient measure cannot have these values
on all cylinder partitions.
\end{proof}

A logit parametrization does not remove this obstruction. If every
independent coordinate has the common infinitesimal logit $t$, then
$\sigma(t)=1/2+t/4+O(t^3)$. Its exponent-one cylinder coefficients are
one quarter of those in \eqref{fsp:eq:coinrow}, and their variations are
still unbounded. Likewise, with $u=\exp(-L)$ for positive infinite $L$,
a shared success probability $\sigma(-L)=u/(1+u)$ has coefficient one
at the first $u$ scale in every first-success event. The first proof
applies in a Hahn chart in $u$. Specifying infinite or infinitesimal
logits therefore does not supply a missing infinite-addition rule.

The unperturbed fair coin has an ordinary real product law, hence a
coefficientwise law at exponent zero. An infinitesimal perturbation at
every coordinate is not automatically a harmless perturbation of the
infinite product: its first coefficient already has unbounded variation.
This is a precise failure of extension, not a claim that every
non-Archimedean Kolmogorov-type theorem fails. Existing theorems with
other hypotheses and addition rules address different questions
\cite{bhw,brickhill,repomeasures}.

\section{Regular all-subsets laws and extension inside the surreals}
\label{fsp:sec:extension}

\subsection{A relative ordered-field embedding lemma}

To distinguish an obstruction in one workspace from absolute
inconsistency, we need an explicit enlargement principle.

\begin{lemma}[Set-sized relative field embedding]\label{fsp:lem:embedding}
Let $K\subseteq\No$ be a set-sized ordered subfield, and let $L/K$ be
a set-sized ordered field extension. There is an ordered-field embedding
$j:L\to\No$ fixing $K$ pointwise.
\end{lemma}
\begin{proof}
Replace $L$ by its ordered real closure. Its elements algebraic over
$K$ form a real closure of $K$, whose unique ordered embedding into
the real closure of $K$ inside $\No$ extends the given identity.
Choose and well order a transcendence basis of $L$ over that real
closure. We extend the embedding successively, keeping the domain and
image real closed after each step.

At a step let $F$ be the current real closed domain and $x$ the next
transcendental basis element. The two sets
\begin{equation}
 \{j(a):a\in F,\ a<x\},\qquad
 \{j(b):b\in F,\ x<b\}
\end{equation}
form a surreal cut. Choose a surreal $y$ strictly between them. It cannot
belong to $j(F)$, since every element of that field lies on one side of
the cut. Since $j(F)$ is real closed, $y$ is transcendental over it.
Sending $x$ to $y$ preserves signs of polynomials: over a real closed
field they factor into linear factors and quadratic factors of constant
sign, and the signs of the linear factors are determined by the cut.
It therefore defines an ordered embedding of $F(x)$, which extends
uniquely to its ordered real closure. At limit stages take unions;
a union of a chain of real closed fields is real closed. Finally $L$
is the real closure of the field generated by its transcendence basis,
so the construction embeds all of $L$. All stages and cuts are sets.
\end{proof}

The map preserves field operations and order. It need not preserve a
previously chosen exponential on $L$, an internal-set structure, or a
specified Hahn summation operation. After embedding we may take the
canonical exponential closure from Lemma~\ref{fsp:lem:workspace}, but that
is not a proof that two different exponential structures agree.

\subsection{A fine-ultrafilter construction}\label{fsp:sub:ultrafilter}

The following is a direct version of the non-Archimedean sampling
construction in the literature \cite{bhw,brickhill}. Let $\Omega$ be
an infinite set and let $I$ be its nonempty finite subsets. For each
finite $F\subseteq\Omega$, the cone
$\{s\in I:F\subseteq s\}$ is nonempty, and finite intersections of
cones are cones. Extend the cone filter to an ultrafilter $\mathcal U$.
This use of the ultrafilter lemma is nonconstructive.

The ultrapower $L=\R^I/\mathcal U$ is an ordered field. Define
\begin{equation}\label{fsp:eq:ultrafilter}
 P(A)=\left[\frac{|A\cap s|}{|s|}\right]_{\mathcal U},\qquad
 H=[|s|]_{\mathcal U},\qquad A\subseteq\Omega.
\end{equation}
These values give an all-subsets regular finitely additive probability.
Indeed normalization and disjoint additivity hold at every snapshot.
For $a\in A$, the cone containing $a$ belongs to $\mathcal U$, so
$P(A)>0$. Fineness also gives
\begin{equation}
 H>n\quad(n\in\N),\qquad P(\{a\})=1/H
 \quad(a\in\Omega).
\end{equation}
Lemma~\ref{fsp:lem:embedding}, applied over $\R$, turns these into surreal
probability values without changing any field identities. This yields
positive equal singleton masses even on an uncountable set.

The addition implicit in \eqref{fsp:eq:ultrafilter} is a snapshot operation:
\begin{equation}
 P(A)=\left[\sum_{a\in A\cap s}\frac1{|s|}\right]_{\mathcal U}.
\end{equation}
It is not a strong Hahn sum of the identical surreal numbers $1/H$.
The finite terms being summed depend on the snapshot before passage to
the quotient. Erasing this dependence would change the construction and
produce the strong-summability obstruction of Section~\ref{fsp:sec:infinity}.

\subsection{Nonuniqueness and limited symmetry}\label{fsp:sub:nonunique}

On $\N_{\ge1}$ one can instead use the initial segments
$s_n=\{1,\ldots,n\}$ and a nonprincipal ultrafilter on the indices.
This remains fine for every fixed finite subset of $\N$. If the
ultrafilter concentrates on even indices, the even numbers have
probability exactly $1/2$. If it concentrates on odd indices, their
probability is
\begin{equation}
 \frac12-\frac{1}{2H},\qquad H=[n].
\end{equation}
For either model the ordered field $\R(H)$ can be identified with
$\R(\omega)$ by $H\mapsto\omega$: both generators are positive
infinite and transcendental over $\R$, so polynomial signs are determined
by their leading terms. Lemma~\ref{fsp:lem:embedding} then embeds the full
field over that identification. The two resulting surreal probabilities
have the same chosen singleton mass $1/\omega$ and the same real
probability of the evens, but different exact probabilities of the evens.
Thus even fixing the infinitesimal unit need not select a unique law.

Equal singleton probabilities imply invariance under permutations with
finite support, by finite additivity. They do not imply invariance
under arbitrary transformations. In particular, a shift-invariant law
on all subsets of $\N$ would satisfy
$P(\N\setminus\{1\})=P(\N)$, contradicting a positive singleton mass.
More generally, an invariance axiom can conflict with regularity even
when every atom has the same mass. Symmetry requirements must be checked
as separate constraints, not inferred from the word ``uniform.''

\subsection{Regular extension of any regular event algebra}

\begin{theorem}[All-subsets regular extension after field enlargement]
\label{fsp:thm:extension}
Let $\Omega$ be a nonempty set, $\A$ an algebra of its subsets, and
$P_0:\A\to K$ a regular finitely additive probability, where
$K\subseteq\No$ is set sized. There is a set-sized ordered field
extension $L/K$ and a regular finitely additive probability
$P:\mathcal P(\Omega)\to L$ extending $P_0$.
The extension can be realized with values in a set-sized subfield of
$\No$ while preserving every original value in $K$.
\end{theorem}
\begin{proof}
Use the first-order language of ordered fields with constants for the
elements of $K$ and a constant $p_A$ for each subset $A\subseteq\Omega$.
Include the ordered-field diagram of $K$, normalization, disjoint finite
additivity, the axioms $p_A>0$ for every nonempty $A$, and
$p_A=P_0(A)$ for $A\in\A$.

Every finite fragment is satisfiable in $K$. To show this, take the
finite Boolean algebra generated by the finitely many original events
from $\A$ named in the fragment. Each of its nonempty atoms has positive
$P_0$ mass. Refine its atoms by all the other finitely many subsets
mentioned in the fragment. Within a parent atom of mass $b>0$, split
$b$ equally among its finitely many nonempty refined children. The
resulting positive finite atomic law has exactly the prescribed masses
on the original finite Boolean algebra and satisfies every finite
additivity constraint in the fragment. Empty cells have mass zero.
The constants from $K$ retain their given values.

First-order compactness gives a set-sized ordered-field model of all
these axioms. Its copy of $K$ is an ordered subfield, and the values of
the $p_A$ define the desired extension. Restricting to the subfield
generated by those values and $K$ is enough. Finally apply
Lemma~\ref{fsp:lem:embedding} over $K$; its image remains a set-sized field
and all original probabilities are fixed.
\end{proof}

The logical compactness theorem is the standard first-order theorem;
no topological compactness of a surreal simplex is being assumed.
Regularity is necessary if a regular extension must preserve every
original mass: an original nonempty zero-mass event cannot become
positive. The theorem does not require countable additivity, symmetry,
a canonical embedding, or a prescribed bound on the new scales.
Imposing such extra conditions may block the construction.

\begin{corollary}[The two coin models do have regular finite-additive laws]
\label{fsp:cor:coinextension}
The finite cylinder probabilities of either Bernoulli$(t)$ or
Bernoulli$(1/2+t)$ extend to regular finitely additive surreal
probabilities on all subsets of $\{0,1\}^{\N}$, with the original
$t=\omega^{-1}$ unchanged.
\end{corollary}
\begin{proof}
Each nonempty cylinder-algebra event is a nonempty finite union of
positive finite-pattern atoms, so the consistent cylinder law is
regular and finitely additive. Apply Theorem~\ref{fsp:thm:extension}
over its original scalar field.
\end{proof}

Theorems~\ref{fsp:thm:rarecoin} and \ref{fsp:thm:faircoin} prove that these
extensions cannot carry the specified coefficientwise finite-measure
structure. There is no contradiction: the extension theorem promised
a different axiom system. It also does not compute the probability of
an arbitrary tail event. Such probabilities can depend on which model
of the extension axioms has been selected.

\section{An infinite learning model with a rare latent regime}
\label{fsp:sec:latent}

The previous obstructions should not leave the impression that no
interesting infinite stochastic models survive coefficientwise addition.
Mixtures of ordinary path laws provide explicit counterexamples to that
impression, and also expose a limitation of finite Bayesian updating.

Let $\mu_0$ be the ordinary independent Bernoulli$(1/3)$ law on
$\{0,1\}^{\N}$ and $\mu_1$ the independent Bernoulli$(2/3)$ law. Attach
a latent label $H\in\{0,1\}$ and give the joint space the Hahn law
\begin{equation}\label{fsp:eq:latentlaw}
 P=\frac{\delta_0\otimes\mu_0+t\delta_1\otimes\mu_1}{1+t}.
\end{equation}
Theorem~\ref{fsp:thm:hierarchy} makes this a positive coefficientwise
probability on its product sigma-algebra. It has an exact rare regime
$P(H=1)=t/(1+t)>0$, not merely a family of finite-dimensional proposals.
Its unconditional coordinate means and covariances are
\begin{equation}\label{fsp:eq:latentcov}
 \E X_i=\frac{1+2t}{3(1+t)},\qquad
 \Cov(X_i,X_j)=\frac{t}{9(1+t)^2}>0\quad(i\ne j).
\end{equation}
To check the covariance, write the mean as the mixture of $1/3,2/3$
and the cross moment as the same mixture of their squares. Subtraction
leaves the variance of the latent success probability. The process is
exchangeable but not independent. Therefore its infinitesimally shifted
marginals do not contradict Theorem~\ref{fsp:thm:faircoin} or the analogous
coefficientwise independent-product obstructions.

Let $\F_n$ record the first $n$ observed coordinates, and write
$K_n=\sum_{i=1}^nX_i$. Bayes' rule gives the finite conditional
probability of the rare regime:
\begin{equation}\label{fsp:eq:latentposterior}
 M_n=P(H=1\mid\F_n)
 =\frac{t\,2^{2K_n-n}}{1+t\,2^{2K_n-n}}.
\end{equation}
Every value of $M_n$ is infinitesimal for each ordinary $n$: the factor
$2^{2K_n-n}$ is a positive real number. These finite-range variables
form a bounded martingale by the finite tower property.

The ordinary strong laws of large numbers of the two real component measures
\cite{kallenberg} give the Borel event
\begin{equation}
 B=\left\{\lim_{n\to\infty}\frac{K_n}{n}=\frac23\right\},
 \qquad \mu_0(B)=0,\quad\mu_1(B)=1.
\end{equation}
Consequently
\begin{equation}\label{fsp:eq:tailposterior}
 P(B)=\frac{t}{1+t}>0,
 \qquad P(H=1\mid B)=1.
\end{equation}

\begin{proposition}[Finite-update shadows need not recover rare tail conditioning]
\label{fsp:prop:tailfailure}
In model \eqref{fsp:eq:latentlaw}, every finite posterior in
\eqref{fsp:eq:latentposterior} has standard part zero, but observing the tail
event $B$ gives conditional probability one to $H=1$. On $B$, the
finite posteriors do not converge to that conditional value even in the
intrinsic ordered-field topology.
\end{proposition}
\begin{proof}
The finite-posterior assertion was proved above. The component strong
laws of large numbers give \eqref{fsp:eq:tailposterior}. For every ordinary $n$ and every
path, $M_n<1/2$ because it is infinitesimal. Thus on $B$ its distance
from the value one remains greater than $1/2$.
\end{proof}

This also has a precise conditional-expectation formulation. For every
observation event $A$,
$P(H=1,A)=P(B\cap A)$, by the two component almost-sure statements.
Thus $\ind_B$ is a version of the conditional expectation of
$\ind_{\{H=1\}}$ given the full observation sigma-algebra, in the
finite-range integral sense. The martingale $(M_n)$ does not recover
this version on $B$ even at a fixed real resolution. Off $B$, convergence
to zero also fails at infinitesimal resolution: $M_n>t^2$ for every
ordinary $n$ and every path. In fact, adjacent posterior differences
have valuation one and positive real leading magnitude, since the
likelihood-ratio factor is multiplied by either $2$ or $1/2$ at each
step. Thus $|M_{n+1}-M_n|>t^2$ everywhere, so the bounded martingale
is not order-Cauchy along any path.

By contrast, the posterior standard parts are identically zero. Their
ordinary limit differs from $\st\ind_B$ exactly on $B$, whose
probability is positive infinitesimal but whose real shadow probability
is zero. This is fully consistent with the ordinary real martingale
convergence theorem. A surreal analogue retaining all positive
infinitesimal events or all infinitesimal resolutions must specify a
different conclusion or impose additional hypotheses.

The example also distinguishes two uses of ``infinitely much evidence.''
A Borel tail observation is an event in a specified path law. It is not
an unrestricted surreal-indexed sum of finite log-likelihood increments.
The first has a rigorous interpretation here; the second has not been
defined merely by writing an infinite logit.

\section{Standard-part bridges, Loeb measure, and sampling limits}
\label{fsp:sec:bridges}

\subsection{Exactly when the real shadow is countably additive}

\begin{theorem}[A continuity test for the shadow]\label{fsp:thm:shadowcountable}
Let $P$ be a finitely additive $K$-probability on a sigma-algebra.
Then $P_0=\st P$ is countably additive exactly when, for every
$A_n\downarrow\varnothing$ and every positive real $\eps$, there is
an $n$ with $P(A_n)<\eps$.
\end{theorem}
\begin{proof}
A real finitely additive probability is countably additive exactly when it is continuous
from above at the empty set. For completeness, countable additivity
gives this continuity by decomposing $A_1$ into successive differences.
Conversely, for disjoint $B_n$ with union $B$, the remainder
$A_N=B\setminus\bigcup_{n\le N}B_n$ decreases to the empty set, so
finite additivity and continuity give
$P_0(B)=\sum_nP_0(B_n)$.

If $P_0(A_n)\to0$, choose $n$ with $P_0(A_n)<\eps/2$.
The infinitesimal difference between $P(A_n)$ and its standard part is
less than $\eps/2$, so $P(A_n)<\eps$.
Conversely, the displayed condition implies that the decreasing real
numbers $P_0(A_n)$ have infimum at most every positive real $\eps$,
hence zero. This is the required continuity.
\end{proof}

The standard part of the normalized hierarchy is the ordinary measure
$\mu_0$ and satisfies this test. The all-subsets uniform law on $\N$
from \eqref{fsp:eq:ultrafilter} does not: singleton shadows are zero while
the shadow of their union is one. More explicitly, for
$A_n=\{n+1,n+2,\ldots\}$, one has
$P(A_n)=1-n/H$ and $P_0(A_n)=1$. Standard part is always a finite-additive
homomorphism of probability values; countable additivity is an extra
property, not an automatic consequence of the existence of standard part.

\subsection{The Loeb bridge requires internal structure}

Loeb's construction starts with an internal finitely additive probability
on an internal algebra in a sufficiently saturated nonstandard setting
\cite{loeb}. Its limited values have real standard parts. On the internal
algebra these define a real premeasure, which extends to a countably
additive measure on the generated sigma-algebra and then to its completion.
This real completed measure is the Loeb measure.

The role of saturation can be seen directly in the finite-probability
case. For a decreasing sequence of internal sets with empty intersection,
countable saturation forces some member to be empty: otherwise their
finite intersections would all be nonempty. In particular, if countably
many disjoint internal sets have an internal union, only finitely many
can be needed to cover that union. Standard-part finite additivity
therefore supplies the premeasure property on the internal algebra.
The ordinary measure-extension theorem then supplies the sigma-algebra
and its countable additivity. These steps concern internal \emph{sets}
as well as non-Archimedean numbers.

An ordered-field embedding of the underlying scalar field into $\No$
preserves standard parts because it fixes $\R$ and preserves order.
Thus one may carry an already specified internal model and its Loeb
construction along with such a scalar embedding. The embedding alone,
however, does not manufacture internal events, saturation, hyperfinite
sets, or transferred sums. A bare surreal $H>\N$ is not a substitute
for this package of structure.

\subsection{A Poisson shadow at reciprocal-probability sample size}\label{fsp:sub:poisson}

Here is a concrete sampling limit with all semantics supplied. Fix a
positive real $\lambda$. For each ordinary integer $N>\lambda$, perform
$N$ independent Bernoulli trials with real success probability
$p_N=\lambda/N$. Let $S_N$ be the count and pass along a nonprincipal
ultrafilter on these indices to the corresponding internal experiments.
Write $H=[N]$ for the hyperinteger sample size. A single trial then has
probability $\lambda/H$.

For a fixed ordinary integer $k\ge0$, the internal event that the count
is $k$ has probability represented by
\begin{equation}\label{fsp:eq:poisson}
 b_k=\left[
 \binom Nk\left(\frac\lambda N\right)^k
           \left(1-\frac\lambda N\right)^{N-k}
 \right],
 \qquad
 \st b_k=e^{-\lambda}\frac{\lambda^k}{k!}.
\end{equation}
For $N<k$ one assigns the probability zero, which affects only finitely
many indices. To verify the standard part, factor the binomial term as
\begin{equation}
 \frac{\lambda^k}{k!}
 \prod_{r=0}^{k-1}\left(1-\frac rN\right)
 \left(1-\frac\lambda N\right)^{N-k}.
\end{equation}
The finite product tends to one in ordinary real analysis. The logarithm
of the last factor is
$(N-k)(-\lambda/N+O(N^{-2}))\to-\lambda$.
This proves the ordinary real limit, and every nonprincipal ultrapower
of a convergent bounded real sequence has that standard part.

In a saturated internal realization, the Loeb probability of the count
being an ordinary finite integer is one: the disjoint events with fixed
counts $k$ have the real Poisson masses in \eqref{fsp:eq:poisson}, whose
ordinary sum is one. Thus the Poisson limit is recovered at a sample
size comparable to the reciprocal of an infinitesimal success probability.
This supplies the missing meaning of the informal scale calculation
$n\asymp1/p$ in a \emph{chosen nonstandard experiment}.

After an ordered-field embedding into $\No$, the probabilities and
standard parts in \eqref{fsp:eq:poisson} remain valid. But an internal power
such as $[(1-\lambda/N)^N]$ must not be silently identified with the
canonical surreal expression $\exp(H\log(1-\lambda/H))$ under a
field-only embedding. The latter uses extra exponential structure.
Equality requires an exponential-compatible representation, not just
preservation of field operations. The fixed-$k$ probability limits
above do not require that extra claim.

\section{An implementation and formalization architecture}
\label{fsp:sec:implementation}

\subsection{Choose a representable scalar fragment}

General surreal numbers are not finite data structures, and equality,
leading-term extraction, or arbitrary infinite-support manipulation is
not provided by a formal normal-form notation alone. A practical engine
should start with a fragment such as rational functions in a positive
infinitesimal, finite multiscale rational expressions, or explicitly
supported Hahn series with a certified leading term. Canonical logarithms
and exponentials may require a larger expression language than the
probability fragment.

For a finite normalized model, a useful representation stores positive
unnormalized weights rather than rounded probabilities. Normalization,
conditioning, and Bayes then use exact sums and quotients. The minimum-plus
and leading-coefficient layer from Theorem~\ref{fsp:thm:scalebayes} can be
computed first, while residual expressions are retained for ties and
later rare conditioning. A leading-scale answer should not be presented
as an exact full probability.

Theorem~\ref{fsp:thm:stability} supplies a precision contract. An operation
conditioning on mass of valuation $\beta$ and requesting output through
scale $\lambda$ must demand input error beyond $\beta+\lambda$.
An implementation which first truncates all atoms at one fixed absolute
threshold cannot guarantee every subsequent conditional query. For
score-based models, Proposition~\ref{fsp:prop:softbound} gives a different
contract: an additive score error bound controls multiplicative
probability error. This can be preferable when tiny probabilities must
remain distinguishable.

\subsection{Separate proof interfaces}

A formal development should not require the full construction of $\No$
for every finite probability theorem. The following layers isolate the
actual assumptions:
\begin{enumerate}
\item An ordered-field finite-probability layer: weights, finite sums,
conditioning, expectation, inequalities, finite products, and stopping.
\item A residue/valuation layer: standard part, leading coefficients,
conditional shadows, minimum-plus updates, and precision theorems.
\item An exponential ordered-field layer: logit, softmax, entropy,
relative entropy, proper scoring, and the explicit Gibbs identity.
\item Separate infinite modules for strong supports, coefficient measures,
compactness/ultrafilter extensions, and internal/Loeb constructions.
\end{enumerate}
The first three modules should not import an unproved global integration
operator. An interface for infinite sums should name its support or
convergence certificate and state which regroupings it permits.
A field embedding interface should distinguish preservation of order,
exponential, standard part, strong sums, and internal structure. None
of the latter properties should be inferred merely from the first.

The repository has a maintained formalization ledger, but no new Lean
proofs are delivered here. The present package's executable tests are
exact checks in a small rational-function fragment. They are evidence
about those examples and the implementation, not a formal verification
of the general mathematical theorems.

\subsection{What the accompanying verifier checks}\label{fsp:sub:verifier}

The file \texttt{code/verify.py} implements exact arithmetic in $\Q(t)$ with
$t$ ordered as a positive infinitesimal. It checks finite Bayesian
normalization and update order, leading-scale posterior formulas,
conditional shadows, conditional precision examples, finite probability
inequalities, the rare-regime covariance, the hidden-shadow examples,
and the explicit cylinder coefficients used in the nonextension proofs.
The delivered run passed 2,145 assertions under Python 3.13.5, with the
breakdown recorded in the JSON output.
All comparisons use rational polynomial arithmetic, not substitution of
a small floating-point number for $t$.

The output \texttt{data/verification.json} records the actual assertion count
and test families. The program cannot check first-order compactness,
uncountable sets, saturation, canonical surreal exponentiation, or an
infinite measure-extension theorem. Those parts are justified by the
written proofs and explicitly cited foundational inputs, not by finite
experiments. The separation is important: thousands of finite successful
checks would not establish an infinite extension axiom.

On placement the unchanged program was rerun on a copy under Python 3.14.4.
It passed the same 2,145 assertions in the same 19 families, and its output,
written to a separate file, matched \texttt{data/verification.json} in every
field except the recorded Python version.

\section{Scope, limitations, and next mathematical questions}
\label{fsp:sec:scope}

\subsection{What has been built}

There is a usable finite probability and decision theory with genuinely
surreal-valued probabilities. It includes exact conditioning on positive
infinitesimal events, canonical log-odds, finite information inequalities,
and explicit optimization identities. It retains a hierarchy of
conditional beliefs invisible to the unconditional real shadow.
There are also genuine infinite-space models: strong atomic examples,
diffuse coefficientwise hierarchies, a rare latent path regime, and
regular all-subsets finite-additive extensions. The relationships among
these constructions are mathematical statements with different hypotheses,
not an assertion that one is the uniquely correct probability theory.

Several results identify useful boundaries rather than failures of the
whole program. The conditioning precision theorem explains which
approximations are safe. The coin obstructions identify coefficient
variation as an extension barrier. The uncountable-support example shows
that coefficientwise measurability alone does not define a positive
integration theory. The latent-regime example separates finite posterior
shadows from conditioning on a tail observation. These are concrete
requirements for any stronger theory.

\subsection{What is not asserted}\label{fsp:sub:notasserted}

We have not constructed a countably additive probability on all events
using the full fine topology of $\No$; Proposition~\ref{fsp:prop:fine}
explains why ordinary set-indexed convergence is unsuitable there.
We have not proved a universal Carath\'eodory, Radon--Nikodym, Fubini,
martingale convergence, strong law of large numbers, central limit, or stochastic calculus
theorem after unrestricted replacement of $\R$ by $\No$. Each such
theorem needs a specified infinite framework. The Poisson and latent
examples import ordinary component laws or a stated nonstandard model;
they are not substitutes for those missing general hypotheses.

No theorem here assigns a unique answer to every exact-zero conditioning
problem, selects a canonical ultrafilter, or reconciles regularity with
all possible symmetry demands. No statement makes arbitrary surreal
expressions effectively computable. No assertion identifies the Conway
omega-map with arbitrary exponential powers or transfers internal
hyperfinite operations through a field-only embedding. The information
and Gibbs results concern an ordinary finite alphabet, not an
unspecified sum over a proper class or infinitely many states.

\subsection{A concrete research agenda}

One productive question is to characterize positive integration domains
for coefficientwise Hahn probabilities beyond bounded real observables
and countable-support functions. Theorem~\ref{fsp:thm:integrationfailure}
shows that a common well-ordered support and integrable coefficients
are insufficient at uncountable rank. A stronger condition could require
simultaneous almost-everywhere control along support filtrations, with a
precise closure theorem under products and conditioning.

A second question is to formulate convergence of conditional beliefs
that preserves selected infinitesimal regimes without demanding full
fine convergence. The rare latent example is a useful test: a proposed
martingale theorem must explain whether its conclusion is taken modulo
standard-null sets, exact-null sets, coefficient ideals, or a chosen
scale cutoff. These choices yield different, potentially useful notions
of learning.

A third question concerns extension with additional structure. Theorem~
\ref{fsp:thm:extension} guarantees regular all-subsets extension after field
enlargement, but not invariance, a preferred support group, or an
exponential-compatible embedding. Identifying exactly which finite
constraints on these extra structures have simultaneous realizations
would turn a broad existence theorem into a more informative modeling
tool. The symmetry counterexample shows that the answer cannot simply
be ``all natural constraints.''

Finally, a verified computational library can make substantial progress
without solving every foundational problem. The finite ordered-field
core, valuation precision contracts, logit identities, and finite
information inequalities have explicit proofs and limited interfaces.
They form a tractable starting point for formalization and exact symbolic
experiments while the infinite theories remain deliberately separate.

\subsection{Conclusion}

Surreal probability is most useful as a multiscale extension of ordinary
probability, not as a license to ignore the meaning of infinite addition.
Its finite algebra is robust: positive infinitesimal events can be
conditioned upon, infinite logits retain nonzero alternatives, and
information and optimization identities remain exact. Its infinite
analysis is plural: support, coefficient variation, internal structure,
and standard-part continuity decide which familiar theorems survive.
Once those choices are explicit, the resulting theory is mathematically
substantial and operationally precise, while its limitations become
statements that can themselves be proved and tested.

\appendix
\section{Notation and dependency ledger}
\label{fsp:app:ledger}

\begin{center}
\small
\begin{tabularx}{\textwidth}{@{}lX@{}}
\toprule
Symbol & Meaning\\
\midrule
$K$ & Set-sized ordered subfield of $\No$ containing $\R$\\
$E$ & Such a field closed under canonical exponential and positive logarithm\\
$t^\gamma$ & Hahn monomial identified with Conway $\omega^{-\gamma}$\\
$v(x),\lc(x)$ & Least support exponent and its coefficient, for $x\ne0$\\
$\st x$ & Real standard part of a limited surreal\\
$\hsum$ & Strong Hahn sum; not an ordinary convergent real series\\
$P_0$ & Unconditional real shadow $\st P$\\
$C(A\mid B)$ & Real shadow of the exact conditional ratio\\
$\lambda_\gamma$ & $-\log(t^\gamma)$, with canonical logarithm\\
$\sigma,\logit$ & Logistic map and its inverse on $(0,1)$\\
$D(p\Vert q)$ & Finite relative entropy, under the support condition\\
$H$ in ultraproducts & A specified internal sample-size class, not an arbitrary ordinal\\
\bottomrule
\end{tabularx}
\end{center}

The algebraic results in Sections~\ref{fsp:sec:finite}--\ref{fsp:sec:bayes} need
only ordered fields plus the explicitly used residue or Hahn structure.
Sections~\ref{fsp:sec:logits}--\ref{fsp:sec:gibbs} additionally use canonical
exponential and its scalar inequalities. Section~\ref{fsp:sec:hierarchy}
uses ordinary real finite measures and Hahn support multiplication.
Section~\ref{fsp:sec:extension} uses choice, ordered real closures, the
surreal cut property, and first-order compactness.
Section~\ref{fsp:sec:latent} imports two ordinary real product laws and their
strong laws of large numbers. Section~\ref{fsp:sec:bridges} distinguishes ordinary real
measure extension from the extra saturation assumption in the Loeb
construction. No infinite result is deduced solely from the finite
exponential-field transfer principle.

\section{Reproduction and evidence boundaries}\label{fsp:app:build}

The source is standalone and uses an internal bibliography, with no
external images, bibliography processor, or downloaded fonts. In the
directory of this report, build and check with
\begin{verbatim}
latexmk -pdf -interaction=nonstopmode -halt-on-error article.tex
python3 code/verify.py --output rerun-verification.json
\end{verbatim}
When \texttt{latexmk} is unavailable, run \texttt{pdflatex} repeatedly
until cross-references and the table of contents stabilize. The program
writes a file only when \texttt{--output} is given; naming a new file keeps
the delivered record \texttt{data/verification.json} intact. The delivered
build script \texttt{code/build.sh} is kept unchanged. As placed it does not
work: it changes into its own directory \texttt{code/} and builds
\texttt{surreal\_probability.tex}, which is not there. The delivery was checked
for unresolved references and typesetting errors, and the PDF was rendered for
visual inspection; \texttt{data/BUILD\_VALIDATION.json} records that check of
the delivered 35-page PDF, before placement.
A clean typesetting log does not certify mathematical correctness.

The research audit lists the consulted repository snapshot and primary
sources. The repository is an evolving collection of research drafts;
this article's references identify the exact comparison point. The
article's own status is likewise a proof-bearing research exposition,
not a claim of independent peer review or exhaustive literature priority.
The finite verifier and the written infinite proofs provide different
kinds of evidence and should not be conflated.

\begin{thebibliography}{99}

\bibitem{conway}
J. H. Conway, \emph{On Numbers and Games}, first edition, Academic Press,
1976; second edition, A K Peters, 2001.

\bibitem{gonshor}
H. Gonshor, \emph{An Introduction to the Theory of Surreal Numbers},
London Mathematical Society Lecture Note Series 110, Cambridge University
Press, 1986. \url{https://doi.org/10.1017/CBO9780511629143}.

\bibitem{vdde}
L. van den Dries and P. Ehrlich, Fields of surreal numbers and
exponentiation, \emph{Fundamenta Mathematicae} \textbf{167} (2001),
173--188. \url{https://doi.org/10.4064/fm167-2-3}.

\bibitem{vddeerr}
L. van den Dries and P. Ehrlich, Erratum to ``Fields of surreal numbers
and exponentiation'', \emph{Fundamenta Mathematicae} \textbf{168}
(2001), 295--297. The original paper should be read with this correction;
no original ordinal-length estimate is used here.

\bibitem{bournez}
O. Bournez and Q. Guilmant, \emph{Surreal fields stable under exponential
and logarithmic functions}, 2022, arXiv:2201.08199.
\url{https://arxiv.org/abs/2201.08199}.

\bibitem{bhw}
V. Benci, L. Horsten, and S. Wenmackers, Non-Archimedean probability,
\emph{Milan Journal of Mathematics} \textbf{81} (2013), 121--151.
\url{https://doi.org/10.1007/s00032-012-0191-x}.
Preprint: \url{https://arxiv.org/abs/1106.1524}.

\bibitem{halpern}
J. Y. Halpern, Lexicographic probability, conditional probability, and
nonstandard probability, \emph{Games and Economic Behavior}
\textbf{68} (2010), 155--179.
\url{https://arxiv.org/abs/cs/0306106}.
Author manuscript: \url{https://www.cs.cornell.edu/home/halpern/papers/lex.pdf}.

\bibitem{brickhill}
H. Brickhill and L. Horsten, \emph{Popper Functions, Lexicographical
Probability, and Non-Archimedean Probability}, 2016, arXiv:1608.02850.
\url{https://arxiv.org/abs/1608.02850}.
The cited title and version are those of this preprint.

\bibitem{hammond}
P. J. Hammond, Elementary non-Archimedean representations of probability
for decision theory and games, in \emph{Patrick Suppes: Scientific
Philosopher}, Synthese Library 234, 1994, 25--61.
\url{https://doi.org/10.1007/978-94-011-0774-7_2}.

\bibitem{spohn}
W. Spohn, Ordinal conditional functions: A dynamic theory of epistemic
states, in W. L. Harper and B. Skyrms (eds.), \emph{Causation in Decision,
Belief Change, and Statistics}, vol. II, 1988, 105--134.
Author's institutional record:
\url{https://kops.uni-konstanz.de/entities/publication/b98c06e9-d6de-49f8-9379-9fd5fa5f35fe}.

\bibitem{loeb}
P. A. Loeb, Conversion from nonstandard to standard measure spaces and
applications in probability theory, \emph{Transactions of the American
Mathematical Society} \textbf{211} (1975), 113--122.
\url{https://doi.org/10.1090/S0002-9947-1975-0390154-8}.

\bibitem{kallenberg}
O. Kallenberg, \emph{Foundations of Modern Probability}, third edition,
Probability Theory and Stochastic Modelling 99, Springer, 2021.
\url{https://doi.org/10.1007/978-3-030-61871-1}.

\bibitem{repomeasures}
\repo, \emph{Hahn-Valued Measures and Probability: Strong atomicity,
coefficientwise extension, and hidden negative mass}, maintained
research report, repository snapshot \pin.
\url{https://github.com/VladimirReshetnikov/Surreal/tree/d22a5b35d5b3040e870c3cfd6d1c8f7259e094b0/docs/surreal/hahn-valued-measures-and-probability}.

\bibitem{repofound}
\repo, \emph{Foundations}, repository report and shared foundations
conventions at snapshot \pin.
\url{https://github.com/VladimirReshetnikov/Surreal/tree/d22a5b35d5b3040e870c3cfd6d1c8f7259e094b0/docs/foundations-and-computation/foundations}.

\bibitem{repomarkov}
\repo, \emph{Markov Generators at Every Scale}, repository report
at snapshot \pin.
\url{https://github.com/VladimirReshetnikov/Surreal/tree/d22a5b35d5b3040e870c3cfd6d1c8f7259e094b0/docs/surreal/markov-generators-at-every-scale}.

\bibitem{repophys}
\repo, \emph{Surreal Scalars in Theoretical Physics: Black-Hole
Singularities, Asymptotic Structure, and the Limits of Replacing the Scalar
Field}, merged research report of the same repository, compared at commit
\texttt{50cb709}; its directory is
\nolinkurl{docs/physics/surreal-scalars-and-spacetime}. This entry was added
on placement.

\medskip\noindent\emph{Placement note.} The three entries above that carry
the manuscript's pin are its own references. This report now belongs to the
same collection, and Section~\ref{fsp:sec:neighbours} compares it with those
reports as committed at \texttt{50cb709}.

\end{thebibliography}
\end{document}
